% Macro safety net for cross-volume universal-conductor references,
% provided locally so the chapter remains portable
% under standalone extraction. Definitions match
% chapters/theory/universal_conductor_K_platonic.tex.
\providecommand{\BRST}{\mathrm{BRST}}
\providecommand{\BRSTGauged}{\mathrm{BRSTGaugedChirAlg}}

% Regime I: quadratic (Convention~\ref{conv:regime-tags}).

\section{The algebra and shadow archetype}
\label{sec:heisenberg-archetype}

The Heisenberg algebra~$\cH_k$ is the Gaussian archetype. Its
rank-one level~$k$ is exactly the modular characteristic:
\[
\kappa(\cH_k)=k.
\]
For a rank-$d$ diagonal Heisenberg algebra
$\cH_\ell^{\oplus d}$ at common level~$\ell$, the trace over
the $d$ generators gives
\[
\kappa(\cH_\ell^{\oplus d})=d\,\ell .
\]
This is the abelian normalization recorded in the landscape census;
it is not the affine formula with $h^\vee=0$, and it is not produced
by a level shift. In the rank-one formulas below the subscript
$\kappa$ denotes the Heisenberg level itself, so
$\kappa(\cH_\kappa)=\kappa$.

The modular Koszul triple has zero scalar complementarity sum
$\kappa(\cH_k)+\kappa(\cH_k^!)=0$, and the shadow tower lies in
quadrichotomy class~G with depth $r_{\max}=2$
(Theorem~\ref{thm:quadrichotomy},
Chapter~\ref{chap:shadow-quadrichotomy-platonic}). The open/closed
MC element keeps the analytic propagator tensors, but its nonlinear
shadow projections vanish beyond the scalar line; the scalar genus
expansion is controlled by~$\kappa(\cH_k)$ and assembles into the
$\hat{A}$-genus.

The value $\kappa(\cH_k) = k$ is the
$\BRSTGauged$-evaluation of the universal conductor functor
$K(\cA) = -c_{\mathrm{ghost}}(\BRST(\cA))$ specialised to the
abelian-CS BRST resolution: a single $bc$ pair carrying ghost charge
$-2k$ resolves $\cH_k$, so
$K(\cH_k) = -(-2k)/2 = k = \kappa(\cH_k)$
(Corollary~\ref{cor:uc-K-heisenberg},
Chapter~\ref{chap:universal-conductor}). The seed
$\kappa(\cH_k)=k$ is verified against the master table of
Chapter~\ref{ch:landscape-census}, where the Heisenberg collision
residue is $r^{\mathrm{coll}}_{\cH}(z)=k/z$.

\begin{center}
\renewcommand{\arraystretch}{1.3}
\begin{tabular}{ll}
\textbf{Algebra} & Heisenberg~$\cH_\kappa$, rank~$1$,
OPE $J(z)\,J(w) \sim \kappa/(z{-}w)^2$ \\
\textbf{Central charge} & $c = 1$ \\
\textbf{Shadow archetype} & G (Gaussian), $r_{\max} = 2$ \\
\textbf{Modular characteristic} & $\kappa(\cH_\kappa) = \kappa$ \\
\textbf{Koszul dual} & $\cH_\kappa^! = \operatorname{Sym}^{\mathrm{ch}}(V^*)$,
$\kappa(\cH_\kappa^!) = -\kappa$ \\
\textbf{Complementarity} & $\kappa_{\mathrm{Heis}}(\cH_\kappa) + \kappa_{\mathrm{Heis}}(\cH_\kappa^!) = 0$ \textup{(}free-field type~I\textup{)} \\
\textbf{$r$-matrix} & $r^{\mathrm{Heis}}_\kappa(z) = \kappa/z$ (single pole, no Lie bracket) \\
\textbf{Genus expansion} & $F_g = \kappa \cdot \lambda_g^{\mathrm{FP}}$ GF $= (\kappa/\hbar^2)(\hat{A}(i\hbar) - 1)$
\end{tabular}
\end{center}

\noindent
The modular Koszul triple
(Definition~\ref{def:modular-koszul-triple}) of the Heisenberg
algebra is
\begin{equation}\label{eq:heisenberg-triple}
\mathfrak{T}_{\cH}
\;=\;
\bigl(\,\cH_\kappa,\;\;
\operatorname{Sym}^{\mathrm{ch}}(V^*),\;\;
r^{\mathrm{Heis}}_\kappa(z) = \kappa/z\,\bigr).
\end{equation}
The Koszul dual $\cH_\kappa^! =
\operatorname{Sym}^{\mathrm{ch}}(V^*)$ is the curved symmetric
chiral algebra on the dual space (not isomorphic to~$\cH_\kappa$
itself), and the $r$-matrix has a single simple pole
(the $d\log$-absorbed residue of the double-pole OPE).

The five bar-Koszul and closed-sector objects remain separate. For
$A=\cH_k$,
\[
\begin{aligned}
A&=\cH_k,\\
B(A)&=T^c(s^{-1}\bar A),\\
A^{\mathrm{i}}&=H^\bullet B(A),\\
A^!&=\mathbb{D}_{\Ran}(A^{\mathrm{i}})
   =\operatorname{Sym}^{\mathrm{ch}}(V^*)\quad
   \text{with }m_0=-k\omega,\\
Z^{\mathrm{der}}_{\mathrm{ch}}(A)
&=\ChirHoch^\bullet(A,A).
\end{aligned}
\]
The equivalence $\Omega B(A)\simeq A$ is bar-cobar inversion. It is
not the construction of the Koszul dual algebra. The Koszul dual
algebra is the Verdier dual of the Koszul-dual coalgebra, not the
level-reflected Heisenberg algebra~$\cH_{-k}$. The derived centre is
the Hochschild closed-sector object; it is not $A$, $B(A)$, $A^{\mathrm i}$,
or~$A^!$.

The Heisenberg bar complex at genus~$0$
(Chapter~\ref{ch:heisenberg-frame}) is controlled by the logarithmic
propagator $d\log(z-w)$. At genus~$1$ this is replaced by the
simple-pole elliptic kernel
$G_\tau(z)=\zeta_\tau(z)-G_2(\tau)z$, with
$G_2(\tau)=\pi^2E_2(\tau)/3$. At genus~$g\ge2$ the kernel is
$\partial_z\log E(z,w)$ for the prime form. Its local coefficients
are jet tensors depending on the point and the chosen coordinate;
only their scalar contractions descend to modular or tautological
classes. The Heisenberg is the simplest nontrivial instance of
Theorem~D: $\kappa$ determines the scalar genus tower
$F_g=\kappa\lambda_g^{\mathrm{FP}}$.

\section{The Heisenberg open--closed scalar}
\label{sec:heisenberg-theta-oc}

\begin{proposition}[Gaussian shadow termination for
 Heisenberg; \ClaimStatusProvedHere]
\label{prop:heisenberg-gaussian-termination}
\label{prop:heisenberg-eisenstein-master-mc}
For the Heisenberg algebra~$\cH_\kappa$, the shadow
obstruction tower terminates at weight~$2$:
$\Theta^{\leq 2}_{\mathrm{scal}} = \kappa$, all higher shadow
obstructions vanish ($o_r = 0$ for $r \geq 3$), and the scalar
projection of the MC element is
\[
\pi_{\mathrm{scal}}(\Theta_{\cH_\kappa})
=\kappa \cdot \eta \otimes \Lambda .
\]
The full tensor kernel still records the genus-dependent propagator.
\end{proposition}

\begin{proof}
The Heisenberg OPE is purely quadratic:
$\alpha^i(z)\,\alpha^j(w) \sim
 \kappa\,\delta^{ij}/(z-w)^2$.
There is \emph{no simple pole}, so the chiral bracket
$\{-,-\} = 0$ (the bracket extracts the $(z-w)^{-1}$
coefficient, which is absent).

\emph{Step~1: All transferred $L_\infty$ brackets of
degree $\geq 3$ vanish.}
The modular $L_\infty$ operations
$\ell_n^{(g)}\colon \cA^{\otimes n} \to \cA$ for
$n \geq 3$ are built from iterated residues at
codimension-$\geq 2$ boundary strata of~$\overline{C}_n$.
Each such residue involves at least one simple-pole
extraction from a pair of fields. Since the Heisenberg OPE
has no simple pole, every such residue vanishes:
$\ell_n^{(g)} = 0$ for all $n \geq 3$, $g \geq 0$.

\emph{Step~2: Obstruction classes vanish.}
The obstruction class $o_r(\cA) \in
H^2(\operatorname{gr}_F^r \mathfrak{g}/\operatorname{gr}_F^{r+1}\mathfrak{g}, d_2)$
at degree~$r$ is a sum of graph amplitudes with $r$
external legs. For $r \geq 3$, each graph amplitude
contains at least one trivalent (or higher) vertex
contributed by some $\ell_n$ with $n \geq 3$.
By Step~1, all such vertices vanish, so $o_r = 0$.

\emph{Step~3: Tower terminates.}
Since $o_r = 0$ for $r \geq 3$, the MC element
$\Theta^{\leq 2}_{\mathrm{scal}} = \kappa$ extends trivially in the
shadow filtration:
$\Theta^{\leq r}_{\mathrm{scal}}=\kappa$ for all $r\ge2$.
For the Heisenberg family, this says
\[
\pi_{\mathrm{scal}}(\Theta_{\cH_\kappa})
= \kappa \cdot \eta \otimes \Lambda,
\]
with $\eta$ the bar propagator and $\Lambda$ the tautological volume
class. The equality is a scalar-shadow statement, not an assertion
that the full MC tensor has no prime-form or Bergman-kernel
dependence. This is exact Gaussianity
(Remark~\ref{rem:unifying-principle}).
\end{proof}

\begin{remark}[Three-pillar interpretation: Heisenberg]
\label{rem:heisenberg-three-pillar}
\index{Heisenberg!three-pillar interpretation}
In the three-pillar architecture
(\S\ref{sec:concordance-three-pillars}):
(i)~the Heisenberg \v{C}ech complex is a \emph{strict} homotopy chiral
algebra ($F_n = 0$ for $n \geq 3$;
Example~\ref{ex:cech-hca-heisenberg}), so Pillar~A adds no new data;
(ii)~the convolution $sL_\infty$-algebra
$\operatorname{hom}_\alpha(\cC^{\textup{!\textasciigrave}}, \cH)$ has all higher
brackets $\ell_k = 0$ for $k \geq 3$; the strict dg~Lie model
\emph{is} the full $sL_\infty$-algebra, and
$\cH^{\mathrm{sh}} = \mathbb{Q}[\kappa]$ is trivially homotopy
invariant (Theorem~\ref{thm:shadow-homotopy-invariance});
(iii)~the log-FM spaces $\overline{\operatorname{FM}}_n(C|D)$ reduce to
ordinary FM for the Heisenberg (no punctures needed in the abelian case),
and the tropicalization is empty beyond depth~$1$.
The Heisenberg is the unique family where all three pillars are strict.
\end{remark}

\begin{remark}[Chart-class enumeration: Heisenberg]
\label{rem:heisenberg-chart-class-enumeration}
\index{chart-class!Heisenberg}
\emph{Type signature:} (Open quadrant; chart presentation;
Beilinson level~$1$; HP = factorisation dg-cat
$\mathcal{C}^{\mathrm{op}}$ on $(X, D, \tau)$, vacuum $b$,
$A_b = \mathrm{End}_\mathcal{C}(b)$; archetype label invariant
under Morita equivalence on the level-$1\to 0$ fibre).
The chart algebra
\[
A_b^{\cH, k} = \mathrm{End}_{\mathcal{C}^{\mathrm{op}}_X}(b^{\cH,k})
\;\simeq\; \cH_k\otimes_{\mathcal{D}_X}(\,\cdot\,)
\]
on the standard vacuum $b^{\cH,k}$ realises the rank-one Heisenberg
chiral algebra at level~$k$. The five-archetype label is~$\mathsf{G}$
(Gaussian, $r_{\max}=2$); on the rank-$d$ diagonal Heisenberg
the label is unchanged. The class-$\mathsf{G}$ chart class is
Morita-rigid: every chart presentation of $\cH_k$ yields the same
Gaussian shadow tower because all higher transferred brackets vanish
identically on the convolution $sL_\infty$-algebra. The Wakimoto
embedding $\cH_k\hookrightarrow V_k(\fg)$ at $\fg=\fgl_1$ identifies
the rank-one Heisenberg chart with the rank-one Cartan oscillator
sub-chart of any affine Kac--Moody chart, so class~$\mathsf{G}$ is
chart-independent on the abelian sub-locus. Within Vol~III the
matching $\Phi_d^{\mathrm{FA}}$-evaluation at the same vacuum yields
the holographic CY companion (Vol~III, level~$1\leftrightarrow 1$),
specialising at $d=2$ to the curved abelian
$\beta\gamma$-system input on $T^*\Sigma$.
\end{remark}

\begin{remark}[Theorem A in four lanes: Heisenberg]
\label{rem:heisenberg-theorem-A-four-lanes}
\index{Theorem A!four lanes!Heisenberg}
On the chart algebra $A_b^{\cH,k}=\cH_k$, Theorem~A
(Theorem~\ref{thm:A-infinity-2} in
Chapter~\ref{chap:theorem-A-infinity-2}) appears in the four lanes
$\mathrm{L}1,\mathrm{L}2,\mathrm{L}3,\mathrm{L}4$:
\begin{itemize}
\item \emph{$\mathrm{L}1$ (chain level, $\mathrm{Ch}(\mathrm{Vect})$).}
The bar--cobar adjunction is realised by the Gaussian Koszul complex
$\bar{B}^{\mathrm{ch}}(\cH_k)\simeq
\mathrm{Sym}^c(V^*[-1])$ with curvature $m_0=-k\,\omega$
(Theorem~\ref{thm:heisenberg-koszul-dual-early}); the witness chain
homotopy is the trivial Gaussian one $h_G=0$ (no transferred higher
bracket).
\item \emph{$\mathrm{L}2$ (Quillen-equivalence).}
The Vallette bar--cobar Quillen equivalence
\cite[Theorem~2.1]{Val16} applies in the Francis--Gaitsgory
factorisation ambient on the Koszul locus
(here the entire chart algebra is Koszul);
the Hinich~\cite{Hinich2003} symmetric-monoidal promotion is
unobstructed because $\cH_k$ is exactly Gaussian.
\item \emph{$\mathrm{L}3$ ($(\infty,1)$-categorical).}
The $(\infty,1)$-localisation gives the adjoint equivalence between
the conilpotent-cofree coalgebra $\Bbarch_X(\cH_k)$ and the
Verdier-dual symmetric chiral algebra $\cH_k^! = \mathrm{Sym}^{\mathrm{ch}}(V^*)$.
\item \emph{$\mathrm{L}4$ ($(\infty,2)$-properad).}
Theorem~$A^{\infty,2}$ (Theorem~\ref{thm:A-infinity-2})
inscribes the $(\infty,2)$-categorical adjoint equivalence on the
properadic envelope; on the Gaussian sub-properad the equivalence is
strictly Quillen, with no higher properadic correction
(class-$\mathsf{G}$ formality).
\end{itemize}
The four lanes coincide on the entire chart locus of $\cH_k$
because $\cH_k$ is Koszul without resonance.
\end{remark}

\begin{remark}[Five-$\kappa$ stratification: Heisenberg]
\label{rem:heisenberg-five-kappa-stratification}
\index{kappa stratification@$\kappa$-stratification!Heisenberg row G}
The five $\kappa$-measurements
$\{\kappa_{\mathrm{cat}}, \kappa^{\mathrm{Hodge}}_{\mathrm{ch}},
\kappa^{\mathrm{Heis}}_{\mathrm{ch}}, \kappa_{\mathrm{BKM}},
\kappa_{\mathrm{fiber}}\}$ on the row $\mathsf{G}$ of the
$5\times 5$ stratification matrix at the rank-one Heisenberg
$\cH_k$ evaluate as
\[
\kappa_{\mathrm{cat}}^{\cH_k}=k,\qquad
\kappa^{\mathrm{Hodge}}_{\mathrm{ch}}(\cH_k)=k,\qquad
\kappa^{\mathrm{Heis}}_{\mathrm{ch}}(\cH_k)=k,\qquad
\kappa_{\mathrm{BKM}}=0,\qquad
\kappa_{\mathrm{fiber}}=0.
\]
The first three coincide because the Heisenberg is the
defining algebra of the chiral Heisenberg projection (the
$\kappa^{\mathrm{Heis}}_{\mathrm{ch}}$ functional is normalised to
agree with the chart scalar on $\cH_k$ itself);
$\kappa_{\mathrm{BKM}}$ vanishes because the Heisenberg is not a
Borcherds--Kac--Moody algebra; $\kappa_{\mathrm{fiber}}$ vanishes
because no fibre datum is present on the abelian chart.
\emph{Type signature:} (Open quadrant; chart-level scalar tuple;
Beilinson level~$5$; HP = five distinct measurements named;
collapse pattern $r_{\max}=2$). five-invariant discipline: the five
$\kappa$-measurements are not a single invariant; their coincidence
on $\cH_k$ is a class-$\mathsf{G}$ collapse, not a definition.
\end{remark}

\begin{remark}[Four-level structure: Heisenberg]
\label{rem:heisenberg-four-level}
The Heisenberg algebra illustrates the four levels of the modular
engine:
(i)~the free boson on a curve (geometric object);
(ii)~the Heisenberg Lie algebra with bar complex
(strict model $\Convstr$);
(iii)~the $L_\infty$-algebra $\Definfmod(\cH)$ is
\emph{trivially formal}: $\ell_n = 0$ for $n \geq 3$, so
the strict model is already the homotopy-invariant object
(Corollary~\ref{cor:strictification-comparison});
(iv)~Gaussian shadow archetype: the extension tower terminates
at $r = 2$ (shadow weight $\leq 2$).
In the arithmetic decomposition of
Theorem~\textup{\ref{thm:depth-decomposition}} (Chapter~\textup{\ref{chap:arithmetic-shadows}}),
$d_{\mathrm{arith}} = 1$ (a single constituent centre from $\zeta(s)$)
and $d_{\mathrm{alg}} = 0$, giving total depth $d = 2$; the
Gaussian class is precisely the one-centre stratum.
\end{remark}

\begin{remark}[Mode--Bergman correspondence and the sewing envelope]%
\label{rem:heisenberg-mode-bergman}%
\index{Heisenberg!mode--Bergman correspondence}%
\index{Bergman space!Heisenberg identification}%
On the algebraic Heisenberg Fock core, define the mode--Bergman map
\[
\Theta(a_{-n-1}\mathbf{1})
\;=\;
\sqrt{n{+}1}\; z^n
\;\in\; A^2(D),
\]
extended multiplicatively by symmetrization:
$\Theta(a_{-n_1-1} \cdots a_{-n_k-1}\mathbf{1})
= \operatorname{Sym}(\Theta(a_{-n_1-1}\mathbf{1}) \otimes
\cdots \otimes \Theta(a_{-n_k-1}\mathbf{1}))$.
This gives a dense algebraic map
$\Theta \colon \text{Fock core} \to
\operatorname{Sym} A^2(D)$
from the Heisenberg Fock space to the symmetric algebra of the
Bergman space of the disk.
Moriwaki~\cite{Moriwaki26b} identifies a conformally flat 2-disk
algebra on $\operatorname{Sym} A^2(D)$ with the ind-Hilbert completion
of the affine Heisenberg vertex algebra. The conditional sewing
theorem below uses this identification as the analytic completion
input for the algebraic Heisenberg VOA.
\end{remark}

\begin{theorem}[Heisenberg sewing theorem;
\ClaimStatusConditional]\label{thm:heisenberg-sewing-vol1}%
\index{Heisenberg!sewing theorem}%
Under the mode--Bergman correspondence~$\Theta$
\textup{(}Remark~\textup{\ref{rem:heisenberg-mode-bergman}}\textup{)}:
\begin{enumerate}[label=\textup{(\roman*)}]
\item \textup{(\ClaimStatusProvedElsewhere, Moriwaki~\cite{Moriwaki26b})}
 The sewing envelope of the algebraic Heisenberg VOA is exactly
 $\operatorname{Sym} A^2(D)$.
\item \textup{(\ClaimStatusProvedHere)}
 The completed bar differential is the closure of the Gaussian
 collision operator on Bergman tensors.
\item \textup{(\ClaimStatusProvedHere)}
 Pair-of-pants sewing is the second quantization of a
 one-particle Bergman restriction/composition operator.
\item \textup{(\ClaimStatusProvedHere)}
 Every closed amplitude is a Fredholm determinant on the
 one-particle Bergman space.
\end{enumerate}
\end{theorem}

\begin{proof}
Clause~(i) is the ind-Hilbert identification of
Moriwaki~\cite{Moriwaki26b}, which exhibits a conformally flat 2-disk
algebra structure on~$\operatorname{Sym} A^2(D)$ isomorphic to the
Heisenberg completion.
Clause~(ii) follows from the analytic-algebraic comparison theorem
(Theorem~\ref{thm:analytic-algebraic-comparison}), which identifies
the algebraic bar differential with the closure of the Gaussian
collision operator upon passage to the Bergman completion.
Clauses~(iii) and~(iv) follow from the one-particle sewing reduction
(Theorem~\ref{thm:heisenberg-one-particle-sewing}): the Wick theorem
reduces the full Fock-space sewing amplitude to a Fredholm determinant
$\det(1 - K)$ on the one-particle Bergman space, where~$K$ is the
composition operator encoding pair-of-pants gluing.
The Hilbert--Schmidt condition
(Proposition~\ref{prop:hs-trace-class}) ensures that~$K$ is
trace-class, so the determinant converges absolutely.
\end{proof}

\begin{remark}[Analytic sewing input]\label{rem:heisenberg-sewing-input}
The analytic input is Moriwaki's identification
of the completed Heisenberg disk algebra with
$\operatorname{Sym}A^2(D)$~\cite{Moriwaki26b}. Once this completion
is fixed, the algebraic bar differential, pair-of-pants sewing, and
closed amplitudes are the Gaussian collision operator, second
quantization of one-particle Bergman restriction, and the resulting
Fredholm determinant.
\end{remark}

\section{Collision residue, scalar complementarity, and genus tower}
\label{sec:heisenberg-projections}

The scalar projection
$\pi_{\mathrm{scal}}(\Theta_{\cH_\kappa})
= \kappa \cdot \eta \otimes \Lambda$
is Gaussian (Proposition~\ref{prop:heisenberg-gaussian-termination}).
The five standard projections follow.

\subsection{The collision \texorpdfstring{$r$}{r}-matrix}
\label{subsec:heisenberg-r-matrix}
\index{Heisenberg algebra!r-matrix@$r$-matrix}

\begin{proposition}[Heisenberg $r$-matrix; \ClaimStatusProvedHere]
\label{prop:heisenberg-r-matrix}
The collision residue of the MC element is:
\begin{equation}\label{eq:heisenberg-r-matrix}
r^{\mathrm{Heis}}_\kappa(z)
\;=\; \operatorname{Res}^{\mathrm{coll}}_{0,2}(\Theta_{\cH_\kappa})
\;=\; \frac{\kappa}{z}.
\end{equation}
This is a single pole: the Heisenberg $r$-matrix has no
higher-pole terms.
\end{proposition}

\begin{proof}
The OPE $J(z)\,J(w) \sim \kappa/(z{-}w)^2$ has a double pole.
The bar construction extracts the residue along
$d\log(z{-}w) = dz/(z{-}w)$, absorbing one power of
$(z{-}w)$ (Remark~\ref{rem:heisenberg-bar-absorbs-pole} below).
Put $t=z-w$. The singular OPE measure is
$\kappa\,t^{-2}\,dt$. Since $dt=t\,d\log t$, the same singularity
written in the $d\log$ measure has coefficient
$\kappa\,t^{-1}$. Thus
\[
r^{\mathrm{coll}}_{\cH}(t)=\frac{\kappa}{t}.
\]
There is no $(z{-}w)^{-1}$ term in the OPE (no Lie bracket), so
$r^{\mathrm{coll}}_{\cH}(t)$ has no bracket contribution:
$r^{\mathrm{coll}}_{\cH}(t)=\kappa/t$ exactly.
\end{proof}

\begin{remark}[Pole absorption and the pre-dualisation distinction]
\label{rem:heisenberg-bar-absorbs-pole}
\index{pole absorption!Heisenberg example}
The collision $r$-matrix
$r^{\mathrm{Heis}}_\kappa(z)
= \operatorname{Res}^{\mathrm{coll}}_{0,2}(\Theta_{\cH_\kappa})$
lives one pole order below the OPE
(the bar kernel absorbs one pole): the OPE has a
double pole $\kappa/z^2$; the collision residue has a single pole
$\kappa/z$. The $d\log$ extraction sends $z^{-n}$ to
$z^{-(n-1)}$. For a general chiral algebra with OPE
poles at orders $n_1 > n_2 > \ldots$, the collision residue has
poles at orders $n_1-1 > n_2-1 > \ldots$.

The collision residue
$r^{\mathrm{Heis}}_\kappa(z)=\kappa/z$ is a post-extraction object:
it is the level-prefixed $r$-matrix of the bar complex after
$d\log$~absorption.
The pre-extraction singular OPE
$J(z)\,J(w) \sim \kappa/(z{-}w)^2$ has one higher pole order.
It is not the Koszul dual algebra and not the level-reflected
Heisenberg algebra.
\end{remark}

\subsection{Complementarity}
\label{subsec:heisenberg-complementarity}
\index{Heisenberg algebra!complementarity}

\begin{proposition}[Heisenberg complementarity; \ClaimStatusConditional]
\label{prop:heisenberg-complementarity}
The Heisenberg algebra and its Koszul dual satisfy exact
complementarity:
\begin{equation}\label{eq:heisenberg-complementarity}
\text{(free-field type~I):}\qquad \kappa_{\mathrm{Heis}}(\cH_\kappa) + \kappa_{\mathrm{Heis}}(\cH_\kappa^!) \;=\; 0 \qquad \text{(not the Virasoro $13$).}
\end{equation}
The dual algebra $\cH_\kappa^! =
\operatorname{Sym}^{\mathrm{ch}}(V^*)$ is the symmetric chiral
algebra on the dual space, with curvature
$m_0=-\kappa\cdot\omega$. The source vertex algebra
$\cH_\kappa$ is not replaced by $\cH_{-\kappa}$; its genus-$g$
bar curvature is the scalar class $\kappa\lambda_g$, while the
Verdier-dual algebra carries the opposite scalar characteristic.
\end{proposition}

\begin{proof}
The Koszul dual $\cH_\kappa^! = \operatorname{Sym}^{\mathrm{ch}}(V^*)$
is the symmetric chiral algebra on the dual space
(Theorem~\ref{thm:heisenberg-bar-complex-genus0}, \S\ref{thm:frame-heisenberg-koszul-dual});
it has modular characteristic $\kappa(\cH_\kappa^!) = -\kappa$, giving
the free-field pairing
$\kappa_{\mathrm{Heis}}+\kappa_{\mathrm{Heis}}'=0$. This is the
abelian case of Theorem~\ref{thm:modular-characteristic}(iv).
The identity is scalar complementarity; it is not a statement that
$\Theta_{\cH_\kappa^!}=-\Theta_{\cH_\kappa}$ as full MC elements.
\end{proof}

\begin{remark}[Quantum complementarity]
\label{rem:heisenberg-quantum-complementarity}
At genus~$g$, the complementarity theorem
(Theorem~\ref{thm:quantum-complementarity}) gives:
\[
Q_g(\cH_\kappa) + Q_g(\cH_\kappa^!)
\;=\; H^*(\overline{\mathcal{M}}_g,\, Z(\cH_\kappa)).
\]
Here $Z(\cH_\kappa)$ is the ordinary centre local system appearing in
the scalar complementarity statement; it is not the bulk
$Z^{\mathrm{der}}_{\mathrm{ch}}(\cH_\kappa)
=\ChirHoch^\bullet(\cH_\kappa,\cH_\kappa)$.
The Heisenberg is the archetype where the decomposition is
maximally symmetric: both sides contribute equally with
$F_g(\cH_\kappa) = \kappa \cdot \lambda_g^{\mathrm{FP}}$ and
$F_g(\cH_\kappa^!) = -\kappa \cdot \lambda_g^{\mathrm{FP}}$,
and the sum $F_g + F_g' = 0$ reflects the exact cancellation
of the scalar projections.
\end{remark}

\subsection{Genus expansion with Eisenstein series}
\label{sec:heisenberg-genus-expansion-eisenstein}

\subsubsection{Eisenstein series and modular forms}
\label{subsec:eisenstein-recollection}

At genus~$1$ two normalizations occur. The normalized Eisenstein
series $E_{2m}$ has constant term~$1$. The lattice Eisenstein series
is
\[
G_{2m}(\tau)
=\sum_{(a,b)\ne(0,0)}(a+b\tau)^{-2m}
=2\zeta(2m)E_{2m}(\tau)\qquad (m\ge2),
\]
while the regularized weight-$2$ value is
$G_2(\tau)=\pi^2E_2(\tau)/3$. The simple-pole Heisenberg kernel is
the primitive of the double-pole current kernel:
\[
G_\tau(z)=\frac1z-G_2(\tau)z-G_4(\tau)z^3-G_6(\tau)z^5-\cdots .
\]
Thus $E_2$ enters first through $G_2$ and is quasi-modular; the
higher $G_{2m}$ are scalar multiples of the modular forms $E_{2m}$.

\begin{definition}[Eisenstein series]\label{def:eisenstein-series}
\index{Eisenstein series|textbf}
For even $k \geq 2$, the normalized weight-$k$ Eisenstein series is:
\begin{equation}
E_k(\tau) = 1 - \frac{2k}{B_k} \sum_{n=1}^{\infty} \sigma_{k-1}(n) q^n
\end{equation}
where $q = e^{2\pi i \tau}$ is the nome, $B_k$ are the Bernoulli numbers, and $\sigma_{k-1}(n) = \sum_{d|n} d^{k-1}$ is the divisor function.

Explicitly:
\begin{align}
E_2(\tau) &= 1 - 24\sum_{n=1}^{\infty} \sigma_1(n) q^n
= 1 - 24q - 72q^2 - 96q^3 - 168q^4 - \cdots \label{eq:E2-expansion} \\
E_4(\tau) &= 1 + 240\sum_{n=1}^{\infty} \sigma_3(n) q^n
= 1 + 240q + 2160q^2 + 6720q^3 + \cdots \label{eq:E4-expansion} \\
E_6(\tau) &= 1 - 504\sum_{n=1}^{\infty} \sigma_5(n) q^n
= 1 - 504q - 16632q^2 - 122976q^3 - \cdots \label{eq:E6-expansion}
\end{align}
\end{definition}

\begin{proposition}[Modular transformation laws \cite{Kac}; \ClaimStatusProvedElsewhere]\label{prop:eisenstein-modular}
Under $\tau \mapsto \gamma \cdot \tau = \frac{a\tau + b}{c\tau + d}$ with
$\gamma = \begin{pmatrix} a & b \\ c & d \end{pmatrix} \in SL_2(\mathbb{Z})$:

\emph{Weight~4 and~6 (holomorphic modular).}
\begin{align}
E_4\left(\frac{a\tau+b}{c\tau+d}\right) &= (c\tau + d)^4 E_4(\tau) \\
E_6\left(\frac{a\tau+b}{c\tau+d}\right) &= (c\tau + d)^6 E_6(\tau)
\end{align}

\emph{Weight~2 (quasi-modular).}
\begin{equation}
E_2\left(\frac{a\tau+b}{c\tau+d}\right) = (c\tau + d)^2 E_2(\tau)
+ \frac{6c(c\tau + d)}{\pi i}
\end{equation}

The extra term for $E_2$ is the holomorphic modular anomaly.
\end{proposition}

\begin{definition}[Dedekind eta function]\label{def:eta-function}
\index{Dedekind eta function|textbf}
The Dedekind eta function is:
\begin{equation}
\eta(\tau) = q^{1/24} \prod_{n=1}^{\infty} (1 - q^n)
= q^{1/24}(1 - q - q^2 + q^5 + q^7 - \cdots)
\end{equation}

This is a modular form of weight $1/2$ with a multiplier system:
\begin{equation}
\eta\left(\frac{a\tau+b}{c\tau+d}\right) = \epsilon(\gamma) (c\tau + d)^{1/2} \eta(\tau)
\end{equation}
where $\epsilon(\gamma)$ is a 24th root of unity determined by $\gamma$.
\end{definition}

\subsection{Genus 0: classical Heisenberg (review)}
\label{subsec:heisenberg-genus-zero}

\begin{theorem}[Genus zero correlation functions \cite{FBZ04}; \ClaimStatusProvedElsewhere]\label{thm:heisenberg-genus-zero}
At genus zero ($\mathbb{CP}^1$), the Heisenberg $n$-point function is:
\begin{equation}
\langle a(z_1) a(z_2) \cdots a(z_n) \rangle_0 =
\begin{cases}
0 & n \text{ odd} \\
\kappa^{n/2} \sum_{\text{pairings}} \prod_{(i,j) \in \text{pairing}} \frac{1}{z_i - z_j}
& n \text{ even}
\end{cases}
\end{equation}

For $n=2$:
\[\langle a(z_1) a(z_2) \rangle_0 = \frac{\kappa}{z_1 - z_2}\]

For $n=4$:
\[\langle a(z_1) a(z_2) a(z_3) a(z_4) \rangle_0 = \kappa^2 \left[
\frac{1}{(z_1-z_2)(z_3-z_4)} + \frac{1}{(z_1-z_3)(z_2-z_4)} +
\frac{1}{(z_1-z_4)(z_2-z_3)}
\right]\]

This is Wick's theorem for the chiral boson; no modular forms appear at genus zero.
The propagator $\kappa/(z_i - z_j)$ used here is the \emph{chiral algebra propagator}
(the residue pairing, yielding a simple pole); the current--current correlator
$\langle \partial a(z_1)\, \partial a(z_2)\rangle = \kappa/(z_1-z_2)^2$ is
the $-\partial_{z_1}$ derivative. The all-genus generalization of the
simple-pole propagator is Theorem~\ref{thm:heisenberg-all-genus} below.
\end{theorem}

\subsection{\texorpdfstring{Genus 1: elliptic functions and $E_2$}{Genus 1: elliptic functions and E 2}}
\label{subsec:heisenberg-genus-one-complete}

\begin{theorem}[Genus-1 Heisenberg bar kernels;
\ClaimStatusProvedHere]\label{thm:heisenberg-genus-one-complete}
On an elliptic curve
$E_\tau = \mathbb{C}/(\mathbb{Z} + \tau\mathbb{Z})$, the
simple-pole Heisenberg bar two-point kernel is
\begin{equation}
\langle a(z_1) a(z_2) \rangle^{\mathrm{bar}}_{E_\tau}
= \kappa \cdot G_\tau(z_1 - z_2)
\end{equation}
where $G_\tau(z)$ is the simple-pole elliptic bar kernel, defined
using the Weierstrass $\zeta$-function:
\begin{equation}
G_\tau(z) = \zeta_\tau(z) - \frac{\pi^2 E_2(\tau)}{3}\, z
\end{equation}
where:
\begin{equation}
\zeta_\tau(z) = \frac{1}{z} + \sum_{(m,n) \neq (0,0)} \left[
\frac{1}{z - m - n\tau} + \frac{1}{m+n\tau} + \frac{z}{(m+n\tau)^2}
\right]
\end{equation}
All bar-kernel $n$-point functions are obtained by Wick contraction:
odd correlators vanish, and for $n=2r$,
\begin{equation}
\left\langle \prod_{\nu=1}^{2r}a(z_\nu)\right\rangle^{\mathrm{bar}}_{E_\tau}
=\kappa^r
\sum_{\mathrm{pairings}\ P}
\prod_{\{i,j\}\in P}G_\tau(z_i-z_j).
\end{equation}
The current--current elliptic correlator is the double-pole carrier
$\kappa(-\partial_zG_\tau)(z)$, not the bar kernel itself.
\end{theorem}

\begin{proof}

\emph{Step~1: Weierstrass regularization.}

On the torus $E_\tau$, we need a function with a simple pole at the origin. By Liouville's theorem, no meromorphic function on a compact Riemann surface can have a single simple pole and be doubly periodic. The Weierstrass $\zeta$-function provides the unique regularization of the lattice sum:
\[\zeta_\tau(z) = \frac{1}{z} + \sum_{(m,n) \neq (0,0)} \left[\frac{1}{z - m - n\tau} + \frac{1}{m+n\tau} + \frac{z}{(m+n\tau)^2}\right]\]
This is absolutely convergent and has quasiperiodicity $\zeta(z+1) = \zeta(z) + 2\eta_1$, $\zeta(z+\tau) = \zeta(z) + 2\eta_2$, where $\eta_1 = \pi^2 E_2(\tau)/6$ is the quasi-period (so that $2\eta_1 = \pi^2 E_2(\tau)/3$).

\emph{Step~2: The $E_2$ correction.}

The regularized Green function is:
\[G_\tau(z) = \zeta_\tau(z) - \frac{\pi^2 E_2(\tau)}{3}\, z\]
The subtraction of $2\eta_1 z$ cancels the quasi-period under $z \mapsto z + 1$, making $G_\tau$ periodic in the $1$-direction: $G_\tau(z+1) = G_\tau(z)$. In the $\tau$-direction, $G_\tau(z+\tau) = G_\tau(z) - 2\pi i$ (by the Legendre relation $2\eta_1\tau - 2\eta_2 = 2\pi i$). The residual quasi-periodicity $-2\pi i$ is topological and unavoidable for a function with a single simple pole on a torus.

\emph{Step~3: Laurent expansion.}

Expanding $G_\tau$ near $z=0$ gives the normalized Eisenstein
coefficients
\begin{align}
G_\tau(z)
&= \frac{1}{z} - G_2(\tau)z - G_4(\tau)z^3
-G_6(\tau)z^5-\cdots  \\
&= \frac{1}{z} - \frac{\pi^2 E_2(\tau)}{3} z
-\frac{\pi^4 E_4(\tau)}{45}z^3
-\frac{2\pi^6 E_6(\tau)}{945}z^5-\cdots .
\end{align}

\emph{Step~4: Eisenstein series appearance.}

The coefficient of $z$ in the expansion involves the regularized
lattice Eisenstein series:
\[G_2(\tau) = \sum_{(m,n) \neq (0,0)} \frac{1}{(m+n\tau)^2} = \frac{\pi^2}{3} E_2(\tau)\]

Wick's theorem for the Gaussian Heisenberg field gives the displayed
$n$-point formula. Differentiating the simple-pole kernel gives the
typed current carrier
\[
(-\partial_zG_\tau)(z)=\wp_\tau(z)+G_2(\tau),
\]
so the current OPE keeps its double-pole principal part while the
ordered bar complex uses the primitive simple-pole kernel.
\end{proof}

\begin{remark}[Holomorphic anomaly]\label{rem:E2-holomorphic-anomaly}
\index{holomorphic anomaly!E2 in Heisenberg genus expansion@$E_2$ in Heisenberg genus expansion}
The appearance of $E_2$ introduces the genus-$1$ holomorphic
anomaly. With $z'=z/(c\tau+d)$ and
$\gamma\tau=(a\tau+b)/(c\tau+d)$, the theta-kernel transformation is
\begin{equation}
G_{\gamma\tau}(z')
=(c\tau+d)G_\tau(z)+2\pi i\,c\,z .
\end{equation}
Equivalently,
\[
\langle a(z'_1)a(z'_2)\rangle_{\gamma\tau}
=(c\tau+d)\langle a(z_1)a(z_2)\rangle_\tau
+2\pi i\,c\,\kappa\,z_{12}.
\]
The same anomaly is the inhomogeneous term in
\[
E_2(\gamma\tau)
=(c\tau+d)^2E_2(\tau)
+\frac{6c(c\tau+d)}{\pi i}.
\]
Scalar modular descent replaces $E_2$ by
$\widehat{E}_2=E_2-3/(\pi\operatorname{Im}\tau)$; the loss of
holomorphicity is the genus-$1$ obstruction computed in
Theorem~\ref{thm:heisenberg-obs}.
\end{remark}

\begin{computation}[Partition function at genus 1]\label{comp:heisenberg-partition-g1}
The genus-1 partition function is:
\begin{equation}
Z_{E_\tau}^{\mathcal{H}} = \text{Tr}_{H_{\mathcal{H}}} q^{L_0 - c/24}
= \frac{1}{\eta(\tau)}
\end{equation}

Explicitly:
\begin{align}
Z_{E_\tau}^{\mathcal{H}} &= q^{-1/24} \prod_{n=1}^{\infty} \frac{1}{1 - q^n} \\
&= q^{-1/24}(1 + q + 2q^2 + 3q^3 + 5q^4 + 7q^5 + \cdots)
\end{align}

The coefficients are the partition function $p(n)$ counting partitions of $n$.

Under $\tau \mapsto -1/\tau$, using $\eta(-1/\tau) = \sqrt{-i\tau}\,\eta(\tau)$:
\[Z_{E_{-1/\tau}}^{\mathcal{H}} = \frac{1}{\sqrt{-i\tau}} \cdot Z_{E_\tau}^{\mathcal{H}}\]

since $Z = 1/\eta(\tau)$ transforms contravariantly to $\eta$.
\end{computation}

\subsection{\texorpdfstring{Genus 2: Siegel modular forms $E_4$ and $E_6$}{Genus 2: Siegel modular forms E 4 and E 6}}
\label{subsec:heisenberg-genus-two}

\begin{theorem}[Genus-2 Heisenberg kernel; \ClaimStatusProvedHere]\label{thm:heisenberg-genus-two}
On a genus-2 Riemann surface $\Sigma_2$ with period matrix
$\Omega = \begin{pmatrix} \tau_1 & z \\ z & \tau_2 \end{pmatrix} \in \mathcal{H}_2$,
the Heisenberg two-point function is:
\begin{equation}
\langle a(w_1) a(w_2) \rangle_{\Sigma_2} = \kappa \cdot G_{\Omega}(w_1, w_2)
\end{equation}
where
\[
G_\Omega(w_1,w_2)=\partial_{w_1}\log E(w_1,w_2;\Omega)
\]
and $E$ is the prime form. The Heisenberg algebra is Gaussian:
the two-point function is this kernel with scalar prefactor
$\kappa$, with no perturbative corrections in~$\kappa$. The local
coefficients of $G_\Omega$ are prime-form jet tensors. Their scalar
Siegel modular contractions, when such contractions are formed, lie
in Igusa's genus-2 modular-form ring generated by
$E_4^{(2)},E_6^{(2)},\chi_{10},\chi_{12}$; the local coefficients
themselves are not canonically equal to $E_4^{(2)}$ and
$E_6^{(2)}$.
\end{theorem}

\begin{proof}
Since the Heisenberg algebra is a free (Gaussian) theory, the genus-$g$
two-point function is exactly the propagator, the Green function
$G_\Omega(w_1, w_2) = \partial_{w_1}\log E(w_1, w_2; \Omega)$, with
no perturbative corrections. We construct this explicitly for $g = 2$.

\emph{Step~1: Prime form construction.}

For an odd theta characteristic~$\delta$, the genus-2 prime form is
\[
E(w_1,w_2;\Omega)
=\frac{\theta[\delta](\int_{w_2}^{w_1}\omega,\Omega)}
{h_\delta(w_1)h_\delta(w_2)},
\]
where $h_\delta$ is the corresponding half-differential. This is a
$(-1/2,-1/2)$-form with a simple zero on the diagonal.

\emph{Step~2: Laurent expansion of the genus-2 Green function.}

In a local coordinate $t=w_1-w_2$, the chiral propagator has the
form
\[
G_\Omega(w_1,w_2)=\frac1t+\sum_{m\ge0}c_m(w_2;\Omega)t^m .
\]
The coefficients $c_m$ are local tensors built from the Bergman
projective connection and its covariant derivatives. They depend on
the point and on the local coordinate, so they are not scalar Siegel
modular forms by themselves. This is not a loop correction; the
Heisenberg theory is free and the kernel is exact.

\emph{Step~3: Siegel modular forms.}

For genus~$2$, scalar modular functions obtained from global
contractions of these tensors live in Igusa's ring. The Siegel
Eisenstein series of weight~$k$ is
\[
E_k^{(2)}(\Omega) = \sum_{(C,D)} \det(C\Omega + D)^{-k}
\]
where the sum is over representatives for the bottom rows of
$\mathrm{Sp}_4(\mathbb{Z})$ matrices. Igusa's theorem supplies the
ambient scalar ring
$\mathbb{C}[E_4^{(2)},E_6^{(2)},\chi_{10},\chi_{12}]$ with the
standard relation in even weight. Thus the Heisenberg kernel carries
period-matrix dependence through the prime form, while the scalar
modular projection is governed by the genus-2 Siegel modular ring.
\end{proof}

\begin{remark}[Prime-form jet carrier]\label{rem:genus-two-prime-form-carrier}
The genus-2 kernel has simple-pole principal part and prime-form
jet coefficients. Weight-$4$ and weight-$6$ Siegel Eisenstein series
enter after scalar modular contraction; cusp generators also occur in
the full scalar ring. The Heisenberg theory remains free at every
genus.
\end{remark}

\begin{theorem}[Genus-2 obstruction class for \texorpdfstring{$\mathcal{H}_\kappa$}{H kappa}; \ClaimStatusConditional]\label{thm:heisenberg-genus2-obstruction}
\index{obstruction class!genus-2 Heisenberg}
The scalar genus-2 projection of the Heisenberg bar curvature is:
\begin{equation}\label{eq:heisenberg-genus2-curvature}
(d^{(2)}_{\mathrm{scal}})^2 = \kappa \cdot \lambda_2
\end{equation}
where $\lambda_2 = c_2(\mathbb{E}) \in
H^4(\overline{\mathcal{M}}_2)$ is the second Chern class of the
rank-$2$ Hodge bundle. The obstruction class lands in the Hochschild
bulk:
\begin{equation}
\operatorname{obs}_2(\mathcal{H}_\kappa)
= \kappa \cdot \lambda_2
\in H^4\!\left(\overline{\mathcal{M}}_2,\,
Z^{\mathrm{der}}_{\mathrm{ch}}(\mathcal{H}_\kappa)\right).
\end{equation}
The coefficient is the Heisenberg level, equivalently
$\kappa_{\mathrm{obs}}(\mathcal{H}_\kappa)=\kappa$ in
Definition~\ref{def:genus-g-obstruction}.
\end{theorem}

\begin{proof}
\emph{Step~1: Period matrix and propagator.}
Fix a genus-2 surface $\Sigma_2$ with symplectic homology basis
$(A_1, A_2, B_1, B_2)$ and period matrix
$\Omega = \begin{pmatrix} \tau_1 & z \\ z & \tau_2 \end{pmatrix}
\in \mathfrak{H}_2$. On this local chart let
$\omega_1,\omega_2$ be the normalized abelian differentials
($\int_{A_i}\omega_j=\delta_{ij}$,
$\int_{B_i}\omega_j=\Omega_{ij}$). The genus-2 propagator is:
\begin{equation}\label{eq:genus2-propagator}
K^{(2)}(w_1, w_2 | \Omega) = \partial_{w_1} \log E(w_1, w_2; \Omega)
\end{equation}
where $E(w_1, w_2; \Omega)$ is the prime form on $\Sigma_2$. This has a simple pole at $w_1 = w_2$ with residue~$1$.

\emph{Step~2: B-cycle monodromies.}
The propagator satisfies:
\begin{align}
K^{(2)}(w_1 + \delta_{A_j}, w_2) &= K^{(2)}(w_1, w_2) \label{eq:genus2-A-periodicity} \\
K^{(2)}(w_1 + \delta_{B_j}, w_2) &= K^{(2)}(w_1, w_2) - 2\pi i \cdot \omega_j(w_2) \label{eq:genus2-B-quasi}
\end{align}
At genus~1, there was one B-cycle producing a \emph{constant} shift $-2\pi i$ (equation~\eqref{eq:B-cycle-quasi-periodicity}). At genus~2, there are two B-cycles, each producing a shift by the \emph{function-valued} abelian differential $-2\pi i \cdot \omega_j$. The passage from constant to function-valued monodromies is the key new feature at genus~$\geq 2$.

\emph{Step~3: Curvature from the B-cycle defect matrix.}
The curvature arises from the failure of $B$-cycle single-valuedness,
as in genus~$1$ (Theorem~\ref{thm:heisenberg-obs}). At genus~$2$ the
two $B$-cycle defects form a local frame-dependent vector
$(-2\pi i\,\omega_1,-2\pi i\,\omega_2)$. Under a change of symplectic
basis this vector transforms by the Hodge transition matrix; the
scalar obstruction is therefore not the exterior product of two
globally chosen one-forms. It is the invariant determinant of the
Hodge-curvature matrix:
\begin{equation}\label{eq:genus2-curvature-product}
(d^{(2)}_{\mathrm{scal}})^2
= \kappa \cdot
\det\!\left(\frac{F_{\mathbb E}}{2\pi i}\right),
\end{equation}
where $F_{\mathbb E}$ is the curvature of the Hodge connection on the
rank-$2$ bundle~$\mathbb E$.

\emph{Step~4: Identification with $\lambda_2$.}
Chern--Weil theory identifies
$\det(F_{\mathbb E}/2\pi i)$ with
$c_2(\mathbb E)=\lambda_2$. This globalizes the local monodromy
calculation without choosing global holomorphic one-forms on the
Hodge bundle. Therefore:
\begin{equation}
(d^{(2)}_{\mathrm{scal}})^2
= \kappa \cdot \lambda_2 \in H^4(\overline{\mathcal{M}}_2).
\end{equation}
\end{proof}

\begin{computation}[Genus-2 free energy via universality]\label{comp:heisenberg-F2-universality}
By Theorem~\ref{thm:genus-universality} with $\kappa(\mathcal{H}_\kappa) = \kappa$ and the Faber--Pandharipande $\lambda_g$ formula:
\begin{equation}
F_2(\mathcal{H}_\kappa) = \kappa \cdot \frac{2^3-1}{2^3} \cdot \frac{|B_4|}{4!} = \kappa \cdot \frac{7}{8} \cdot \frac{1/30}{24} = \frac{7\kappa}{5760}
\end{equation}
where $B_4 = -1/30$ is the fourth Bernoulli number.

\emph{Consistency check}: At $\kappa = 1$, this gives
$F_2 = 7/5760 = 0.001215\ldots$, matching the
Faber--Pandharipande scalar formula
(Theorem~\ref{thm:dmvv-agreement}). In addition:
\begin{equation}
\frac{F_2}{F_1} = \frac{7/5760}{1/24} = \frac{7}{240}
\end{equation}
which is independent of $\kappa$, a consequence of the universality factorization $F_g = \kappa \cdot \lambda_g^{\mathrm{FP}}$ \end{computation}

\begin{remark}[Genus-2 vs genus-1: structural differences]\label{rem:genus2-structural}
The genus-2 computation reveals three features invisible at genus~1:
\begin{enumerate}
\item \emph{Multi-dimensional Hodge bundle.} At genus~1, $\mathbb{E}$ is a line bundle and $\lambda_1 = c_1(\mathbb{E})$. At genus~2, $\mathbb{E}$ has rank~2, and $\lambda_2 = c_2(\mathbb{E})$ is a determinantal class involving the interaction of both B-cycle monodromies.
\item \emph{Off-diagonal periods.} The period matrix entry $z = \Omega_{12}$ encodes the interaction between the two handles of $\Sigma_2$. This interaction contributes to $\lambda_2$ through the cross-term in $c_2(\mathbb{E}) = \frac{1}{2}(c_1(\mathbb{E})^2 - 2\operatorname{ch}_2(\mathbb{E}))$.
\item \emph{Function-valued monodromies.} The genus-1 B-cycle monodromy is the constant $-2\pi i$; the genus-2 monodromies $-2\pi i \cdot \omega_j(w)$ are sections of $\omega_{\Sigma_2}$. This explains why $\lambda_2$ lives in $H^4$ (not $H^2$): each B-cycle defect contributes a $2$-form on moduli space, and the two defects combine multiplicatively.
\end{enumerate}
\end{remark}

\begin{remark}[Degeneration behavior]\label{rem:genus2-degeneration}
\index{degeneration!genus-2 to genus-1}
As $\Sigma_2$ degenerates, the obstruction class $\operatorname{obs}_2 = \kappa \cdot \lambda_2$ restricts to the boundary divisors of $\overline{\mathcal{M}}_2$:
\begin{itemize}
\item \emph{Separating degeneration}
 ($\delta_1$: $\Sigma_2 \to E_{\tau_1} \cup E_{\tau_2}$
 via a node):
 $\lambda_2|_{\delta_1}
 = \lambda_1(E_{\tau_1}) \cdot \lambda_1(E_{\tau_2})$,
 the product of genus-1 Hodge classes. The two B-cycle
 monodromies decouple, recovering two independent
 genus-1 obstructions.
\item \emph{Non-separating degeneration}
 ($\delta_0$: $\Sigma_2 \to$ nodal genus-1): the period
 matrix entry $z = \Omega_{12} \to 0$, the Hodge bundle
 acquires a rank-1 sub-bundle, and
 $\lambda_2|_{\delta_0} = 0$. The vanishing reflects
 a non-separating node reducing the
 effective genus without producing a separating
 obstruction.
\end{itemize}
This degeneration behavior is consistent with the recursive
structure of the genus expansion
(Theorem~\ref{thm:heisenberg-all-genus}).
\end{remark}

\begin{computation}[Genus-2 zero-mode determinant]\label{comp:partition-genus-two}
The genus-2 zero-mode determinant is:
\begin{equation}
Z_{0,\Sigma_2}^{\mathcal{H}}
= \det(\operatorname{Im}\Omega)^{-1/2}
\end{equation}
where $\Omega$ is the $2 \times 2$ period matrix of $\Sigma_2$.
This is the \emph{zero-mode contribution} to the full partition
function of a single free boson on a genus-2 surface.
(The genus-1 quantity $1/\eta(\tau)$ in
Computation~\ref{comp:heisenberg-partition-g1} is the
\emph{holomorphic} partition function, including the oscillator
determinant. At genus~$2$, the oscillator determinant contributes
additional factors involving theta functions; the full
non-holomorphic partition function involves
$1/\det(\mathrm{Im}\,\Omega)^{1/2}$ times these oscillator
contributions.)

For $\gamma=\bigl(\begin{smallmatrix}A&B\\ C&D\end{smallmatrix}\bigr)
\in Sp_4(\mathbb Z)$,
\[
\operatorname{Im}(\gamma\Omega)
=(C\Omega+D)^{-T}\operatorname{Im}\Omega\,
(\overline{C\Omega+D})^{-1},
\]
hence
\[
Z_{0,\gamma\Omega}^{\mathcal H}
=|\det(C\Omega+D)|\,Z_{0,\Omega}^{\mathcal H}.
\]
This is a non-holomorphic density. It is not a holomorphic Siegel
modular form of weight $-1/2$, and it is not the oscillator
determinant.
\end{computation}

\subsection{\texorpdfstring{General genus $g$: the expansion}{General genus g: the expansion}}
\label{subsec:heisenberg-general-genus}

\begin{theorem}[Heisenberg at general genus; \ClaimStatusProvedHere]\label{thm:heisenberg-all-genus}
For genus $g$, the Heisenberg two-point function is \emph{exact} (no perturbative expansion):
\begin{equation}
\langle a(z_1) a(z_2) \rangle_{\Sigma_g} = \kappa \cdot G_{\Omega_g}(z_1, z_2)
\end{equation}
where $G_{\Omega_g}$ is the genus-$g$ Green function constructed from the prime form $E(z_1, z_2; \Omega_g)$ and the period matrix $\Omega_g \in \mathcal{H}_g$.

The geometry of $\Sigma_g$ enters through $G_{\Omega_g}$, whose Laurent expansion near $z_1 = z_2$ takes the form:
\begin{align}
g=0: \quad & G(z_1, z_2) = \frac{1}{z_1 - z_2} \quad \text{(chiral propagator)} \\
g=1: \quad & G_\tau(z_1, z_2) = \zeta_\tau(z_1-z_2) - G_2(\tau)(z_1-z_2) \\
g \geq 2: \quad & G_{\Omega_g}(z_1, z_2) = \partial_{z_1}\log E(z_1,z_2;\Omega_g).
\end{align}
For $g=1$ the local coefficients are the lattice Eisenstein series
$G_{2m}(\tau)$. For $g\ge2$ they are prime-form jet tensors; scalar
modular forms appear only after choosing global contractions or scalar
projections. The Heisenberg algebra is free: there are no loop
corrections.
\end{theorem}

\begin{proof}
The Heisenberg algebra is a free (Gaussian) theory: Wick's theorem
reduces all correlation functions to products of two-point functions.
Therefore the genus-$g$ two-point function is exactly the Green function
$G_{\Omega_g}(z_1, z_2) = \partial_{z_1}\log E(z_1, z_2; \Omega_g)$,
the chiral derivative of the prime form on $\Sigma_g$.

\emph{Step~1: Genus expansion.}

The full partition function sums over genera: $Z = \sum_{g \geq 0} \hbar^{2g-2} Z_g$. For the Heisenberg algebra (a free theory), each genus-$g$ contribution is exact; it is \emph{not} a perturbative loop correction. Rather, $Z_g$ encodes the propagator on $\Sigma_g$ through the period matrix $\Omega_g$.

\emph{Step~2: Local coefficients.}

At genus~$1$ the elliptic function is translation-invariant, so the
Taylor coefficients are scalar functions of~$\tau$:
$G_2(\tau),G_4(\tau),\ldots$. At genus~$g\ge2$ there is no
translation-invariant coordinate. A Laurent coefficient of
$\partial\log E$ is a section of a jet bundle over the universal
curve, constructed from the Bergman projective connection and its
covariant derivatives. It becomes a scalar modular form only after a
specified contraction or integration.

\emph{Step~3: Scalar projections.}

The scalar genus tower used by Theorem~D is not the list of local
prime-form coefficients. It is the tautological projection
\[
F_g(\cH_\kappa)=\kappa\,\lambda_g^{\mathrm{FP}}.
\]
For $g=1$ its holomorphic derivative is proportional to
$\kappa E_2(\tau)$, with the anomaly described in
Remark~\ref{rem:E2-holomorphic-anomaly}. For $g=2$, scalar
period-matrix functions obtained from global contractions live in
Igusa's Siegel modular-form ring. No statement here asserts a
universal first coefficient $E_{2g}^{(g)}$ for all $g$.
\end{proof}

\subsection{Kernel coefficients and scalar modular projections}
\label{subsec:modular-weights-all-genus}

\begin{table}[ht]
\centering
\caption{Heisenberg kernel normalizations and scalar modular projections}
\label{tab:modular-weights}
\begin{tabular}{|c|c|c|}
\hline
\textbf{Genus $g$} & \textbf{Kernel coefficient} & \textbf{Scalar projection} \\
\hline
0 & $1/z$ & none \\
\hline
1 & $G_2=\pi^2E_2/3$; $G_{2m}=2\zeta(2m)E_{2m}$ & $E_2\mapsto \widehat{E}_2$ anomaly \\
\hline
2 & prime-form jet tensors & Igusa ring: $E_4^{(2)},E_6^{(2)},\chi_{10},\chi_{12}$ \\
\hline
3 & prime-form jet tensors & scalar Siegel forms after contraction \\
\hline
$g\ge4$ & prime-form jet tensors & no universal leading Eisenstein claim \\
\hline
\end{tabular}
\end{table}

\begin{proposition}[Eisenstein normalization and scalar scope
\cite{Igusa62,Klingen67}; \ClaimStatusProvedElsewhere]\label{prop:modular-weight-formula}
The following statements are the scalar facts used in this example.
\begin{enumerate}[label=\textup{(\roman*)}]
\item At genus~$1$, the simple-pole kernel has coefficients
$G_2=\pi^2E_2/3$ and $G_{2m}=2\zeta(2m)E_{2m}$ for $m\ge2$.
Only $E_2$ is quasi-modular.
\item At genus~$2$, scalar modular contractions of prime-form jet
tensors lie in Igusa's ring generated by
$E_4^{(2)},E_6^{(2)},\chi_{10},\chi_{12}$.
\item Klingen's convergence bound says that the Siegel Eisenstein
series $E_k^{(g)}$ converges absolutely for even $k>g+1$. This is a
statement about the existence of scalar Eisenstein series, not a
statement that the first local coefficient of $\partial\log E$ is
$E_{\lceil g+2\rceil_{\mathrm{even}}}^{(g)}$.
\end{enumerate}
\end{proposition}

\begin{proof}
Part~(i) is the standard Laurent expansion of
$\zeta_\tau(z)-G_2(\tau)z$. Part~(ii) is Igusa's description of the
genus-$2$ scalar Siegel modular-form ring. Part~(iii) is Klingen's
convergence theorem for Siegel Eisenstein series. The distinction
between local jet tensors and scalar modular forms is forced by the
coordinate dependence of the prime-form expansion.
\end{proof}

\begin{remark}[Scope]
The propagator pole order is always one in the bar kernel. The
Eisenstein hierarchy is exact in genus~$1$ because the elliptic curve
is a group. In genus~$g\ge2$, the exact object is the prime-form
kernel; scalar Siegel modular forms enter only after global
projection.
\end{remark}

\subsection{Eta function in partition functions}
\label{subsec:eta-partition-functions}

\begin{theorem}[Partition-function normalizations and determinant line; \ClaimStatusProvedHere]\label{thm:eta-appearance}
Let $\mathcal H_\kappa^{\oplus d}$ be the rank-$d$ Heisenberg vertex
algebra, with level $\kappa$ used to normalize the current OPE. Its
chiral oscillator character is
\begin{equation}
\chi_{\mathcal H_\kappa^{\oplus d}}(\tau)
=\operatorname{Tr}_{\mathsf F_d}\!\left(q^{L_0-d/24}\right)
=\eta(\tau)^{-d}.
\end{equation}
The level $\kappa$ changes the normalization of the current two-point
function and hence the ordered-bar scalar $\kappa(\mathcal H_\kappa)$;
it does not change the number of oscillator modes in the Fock trace.

For the non-chiral rank-one free boson on the flat elliptic curve
$E_\tau$, after fixing the target-volume convention, the modular
factor is
\begin{equation}
Z_{E_\tau}^{\mathrm{bos}}
=C_{\mathrm{vol}}\,
(\operatorname{Im}\tau)^{-1/2}|\eta(\tau)|^{-2}.
\end{equation}
In the standard scalar Laplacian normalization behind the
Kronecker-limit formula,
\begin{equation}
(\det{}'_{\zeta}\Delta_{E_\tau})^{-1/2}
=C_\Delta'\,
(\operatorname{Im}\tau)^{-1}|\eta(\tau)|^{-2}.
\end{equation}
The missing factor $(\operatorname{Im}\tau)^{1/2}$ is the normalized
constant-mode measure. Thus the combined zero-mode and oscillator
factor is
\begin{equation}
(\operatorname{Im}\tau)^{1/2}
(\det{}'_{\zeta}\Delta_{E_\tau})^{-1/2}
=C_\Delta\,
(\operatorname{Im}\tau)^{-1/2}|\eta(\tau)|^{-2},
\end{equation}
with $C_{\mathrm{vol}}$, $C_\Delta'$, and $C_\Delta$ depending only
on the chosen metric and volume normalizations. At genus $g$, the
zero-mode period normalization and the oscillator contribution combine
as the Quillen norm of the determinant line of $\bar\partial$ on
$\Sigma_g$, coupled to the period matrix $\Omega_g$. This combined
analytic datum is not the product
$\prod_i\eta(\Omega_{ii})^{-1}$ except in the separating degeneration
where $\Sigma_g$ breaks into elliptic components.
\end{theorem}

\begin{proof}
The rank-$d$ Heisenberg Fock space is a tensor product of $d$ copies of
the polynomial algebra generated by $a_{-n}$, $n>0$. Each mode
contributes $(1-q^n)^{-d}$, and the vacuum-energy shift gives
\[
\operatorname{Tr}_{\mathsf F_d}\!\left(q^{L_0-d/24}\right)
=q^{-d/24}\prod_{n\ge1}(1-q^n)^{-d}
=\eta(\tau)^{-d}.
\]
The parameter $\kappa$ enters the OPE
$J(z)J(w)\sim \kappa(z-w)^{-2}$; rescaling the current changes this
two-point normalization but not the oscillator counting.

For the non-chiral free boson, the Kronecker-limit formula identifies
the zeta-regularized scalar Laplacian determinant on $E_\tau$ with
the factor
$C(\operatorname{Im}\tau)^2|\eta(\tau)|^4$.  Taking the inverse square
root gives
$C'(\operatorname{Im}\tau)^{-1}|\eta(\tau)|^{-2}$. The normalized
constant-mode integration contributes
$(\operatorname{Im}\tau)^{1/2}$, so the full non-chiral rank-one
modular factor is
$(\operatorname{Im}\tau)^{-1/2}|\eta(\tau)|^{-2}$, up to the displayed
constant. For higher genus the same construction is the Quillen norm
of $\det R\Gamma(\Sigma_g,\mathcal O_{\Sigma_g})$, coupled to the
period matrix. Since $\Omega_g$ has off-diagonal periods, no product
of diagonal eta functions is invariantly defined away from the
separating boundary.
\end{proof}

\subsection{Faber--Pandharipande scalar free energy and DMVV separation}
\label{subsec:physics-comparison}

\begin{theorem}[Faber--Pandharipande scalar free energy; \ClaimStatusConditional]\label{thm:dmvv-agreement}
\index{Mumford formula}
\index{free energy!genus expansion}
For the Heisenberg algebra at ordered-bar scalar
$\kappa(\mathcal H_\kappa)=\kappa$, the scalar genus-$g$ free energy
is
\begin{equation}
F_g(\mathcal H_\kappa)
=\kappa\int_{\overline{\mathcal M}_{g,1}}\psi_1^{2g-2}\lambda_g
=\kappa\,
\frac{2^{2g-1}-1}{2^{2g-1}}\,
\frac{|B_{2g}|}{(2g)!}.
\end{equation}
Here $\psi_1$ is the cotangent line class at the marked point and
$\lambda_g=c_g(\mathbb E)$ is the top Chern class of the Hodge bundle.
The integrand has degree $(2g-2)+g=3g-2=\dim\overline{\mathcal M}_{g,1}$.
For $\kappa=1$ the first values are
\begin{equation}
F_1=\frac1{24},\qquad
F_2=\frac7{5760},\qquad
F_3=\frac{31}{967680}.
\end{equation}
This is the Faber--Pandharipande $\lambda_g$ formula \cite{FP03}.
The Dijkgraaf--Moore--Verlinde--Verlinde formula \cite{DMVV} concerns
elliptic genera of symmetric products and their Hecke expansion. It is
a different construction and is not a proof of the ordered-bar scalar
$\kappa$ or of the Hodge integral above.
\end{theorem}

\begin{proof}
The marked point records the scalar insertion in the topological
gravity projection. Thus the relevant tautological class is
$\psi_1^{2g-2}\lambda_g$, whose degree equals the complex dimension
$3g-2$ of $\overline{\mathcal M}_{g,1}$. Faber--Pandharipande
\cite{FP03} evaluate the resulting Hodge integral:
\[
\int_{\overline{\mathcal M}_{g,1}}\psi_1^{2g-2}\lambda_g
=
\frac{2^{2g-1}-1}{2^{2g-1}}\,
\frac{|B_{2g}|}{(2g)!}.
\]
Multiplication by the ordered-bar scalar $\kappa(\mathcal H_\kappa)$
is the Heisenberg specialization of Theorem~\ref{thm:genus-universality}.
Substituting $B_2=1/6$, $B_4=-1/30$, and $B_6=1/42$ gives the three
values displayed in the statement. The DMVV theorem is not used in
this calculation; its symmetric-product Hecke variables live in a
separate elliptic-genus problem.
\end{proof}

\begin{remark}[Scalar topological-gravity projection]\label{rem:witten-perspective}
In the scalar topological-gravity projection, $F_g$ is the
genus-$g$ free energy and $\lambda_g$ is the Hodge obstruction class.
Bernoulli numbers enter through the Faber--Pandharipande
$\lambda_g$ evaluation, equivalently through the asymptotic expansion
of the free-boson determinant.
\end{remark}

\subsection{Genus carriers}
\label{subsec:heisenberg-dictionary-complete}

\begin{table}[ht]
\centering
\caption{Heisenberg genus carriers}
\label{tab:heisenberg-complete-dictionary}
\begin{tabular}{|l|l|l|}
\hline
\textbf{Genus} & \textbf{Modular forms} & \textbf{Interpretation} \\
\hline
$g=0$ & (none) & rational bar kernel \\
\hline
$g=1$ & $E_2(\tau)$, $\eta(\tau)$ & elliptic kernel and determinant \\
\hline
$g=2$ & $E_4^{(2)},E_6^{(2)},\chi_{10},\chi_{12}$ after scalar projection & prime-form kernel \\
\hline
$g=3$ & scalar Siegel forms after contraction & prime-form kernel \\
\hline
$g \geq 4$ & no universal leading Eisenstein claim & prime-form kernel \\
\hline
\end{tabular}
\end{table}

\subsection{Explicit coefficient checks}
\label{subsec:explicit-coefficients-tables}

\begin{table}[ht]
\centering
\caption{Eisenstein series $q$-expansions (first 10 terms)}
\label{tab:eisenstein-q-expansions}
\begin{tabular}{|l|c|}
\hline
$n$ & $\sigma_1(n)$ \\
\hline
1 & 1 \\
2 & 3 \\
3 & 4 \\
4 & 7 \\
5 & 6 \\
6 & 12 \\
7 & 8 \\
8 & 15 \\
9 & 13 \\
10 & 18 \\
\hline
\end{tabular}
\quad
\begin{tabular}{|c|c|}
\hline
$E_2(\tau)$ & Coefficient \\
\hline
$q^0$ & 1 \\
$q^1$ & $-24$ \\
$q^2$ & $-72$ \\
$q^3$ & $-96$ \\
$q^4$ & $-168$ \\
$q^5$ & $-144$ \\
\hline
\end{tabular}
\end{table}

\begin{table}[ht]
\centering
\caption{Free energy values $F_g = \lambda_g^{\mathrm{FP}}$ at $\kappa = 1$ for genera $1 \le g \le 10$}
\label{tab:free-energy-values}
\begin{tabular}{|c|c|c|}
\hline
$g$ & $F_g$ (exact) & $F_g$ (decimal) \\
\hline
1 & $1/24$ & $4.17 \times 10^{-2}$ \\
2 & $7/5760$ & $1.22 \times 10^{-3}$ \\
3 & $31/967680$ & $3.20 \times 10^{-5}$ \\
4 & $127/154828800$ & $8.20 \times 10^{-7}$ \\
5 & $511/24524881920$ & $2.08 \times 10^{-8}$ \\
6 & $1414477/2678117105664000$ & $5.28 \times 10^{-10}$ \\
7 & $8191/612141052723200$ & $1.34 \times 10^{-11}$ \\
8 & $16931177/49950709902213120000$ & $3.39 \times 10^{-13}$ \\
9 & $5749691557/669659197233029971968000$ & $8.59 \times 10^{-15}$ \\
10 & $91546277357/420928638260761696665600000$ & $2.17 \times 10^{-16}$ \\
\hline
\end{tabular}
\end{table}

\section{Higher-rank Heisenberg and multi-boson expansions}
\label{sec:higher-rank-heisenberg-eisenstein}
\index{Heisenberg algebra!higher rank!Eisenstein expansion}

For the rank-$d$ Heisenberg algebra $\mathcal{H}_\ell^{\oplus d}$
(corresponding to $d$ free bosons at common level $\ell$), the genus
expansion inherits a factored structure from the tensor product.

\subsection{\texorpdfstring{Rank-$d$ bar complex at genus 1}{Rank-d bar complex at genus 1}}

\begin{computation}[Rank-\texorpdfstring{$d$}{d} Heisenberg: genus-1 curvature]
\label{comp:rank-d-heisenberg-genus1}
\index{Heisenberg algebra!rank $d$!curvature}

The rank-$d$ Heisenberg algebra has generators
$a^i(z)$ ($i = 1, \ldots, d$) with OPE
$a^i(z) a^j(w) \sim \ell \delta^{ij}/(z-w)^2$.
The stress tensor is $T = \frac{1}{2\ell} \sum_{i=1}^d
\normord{a^i a^i}$, giving central charge $c = d$.

\emph{Genus-1 two-point function.}
On $E_\tau$, the two-point function is:
\[
\langle a^i(z_1) a^j(z_2) \rangle_{E_\tau}
= \ell \delta^{ij} \cdot \widetilde{G}_\tau(z_1 - z_2)
\]
where $\widetilde{G}_\tau(z) = \wp_\tau(z) + \frac{\pi^2}{3} E_2(\tau)$
is the holomorphic regularized current kernel. It is a quasi-Jacobi
carrier; modular descent replaces $E_2$ by
$\widehat E_2=E_2-3/(\pi\operatorname{Im}\tau)$. The kernel
$\widetilde{G}_\tau$ has a double pole at $z=0$ and equals
$-\partial_z G_\tau(z)$, where $G_\tau$ is the simple-pole bar kernel
of Theorem~\ref{thm:heisenberg-genus-one-complete}. Throughout this
chapter, $G_\tau$ denotes the simple-pole ordered-bar propagator and
$\widetilde{G}_\tau$ denotes its negative derivative, the
double-pole current correlator.

\emph{Curvature from first principles.}
We derive the genus-1 curvature by explicit residue computation,
using the $B$-cycle monodromy mechanism.

\emph{Step~1: Bar differential at genus~$1$.}
On the elliptic curve $E_\tau = \mathbb{C}/(\mathbb{Z}+\tau\mathbb{Z})$,
the degree-$1$ bar element
$a^i(z_1) \boxtimes a^j(z_2) \cdot \eta_{12}^{\mathrm{ell}}$
uses the genus-$1$ propagator
$K^{(1)}(z|\tau) = \theta_1'(z|\tau)/\theta_1(z|\tau)$
(i.e., $\eta_{12}^{\mathrm{ell}} = d\log\theta_1(z_1 - z_2|\tau)$).
The bar differential extracts the OPE by taking the residue
at $z_1 = z_2$:
\begin{equation}\label{eq:heisenberg-bar-diff-genus1}
\dfib(a^i \boxtimes a^j \cdot \eta_{12}^{\mathrm{ell}})
= \operatorname{Coeff}_{d\log(z_1-z_2)}
\!\Bigl[\frac{\ell\delta^{ij}\,d(z_1-z_2)}{(z_1-z_2)^2}
\cdot K^{(1)}(z_1 - z_2|\tau)\Bigr]
\cdot |0\rangle.
\end{equation}
Since the Heisenberg OPE is \emph{purely} double-pole
($a^i(z)a^j(w) \sim \ell\delta^{ij}/(z-w)^2$, with no
simple-pole term), the bar differential on degree-$1$ elements
produces only a vacuum contribution: there is no
current-valued output as in the Kac--Moody case.

\emph{Step~2: $B$-cycle monodromy.}
The propagator $K^{(1)}$ is $A$-periodic but
$B$-quasi-periodic~(\ref{eq:B-cycle-quasi-periodicity}):
\begin{equation}\label{eq:heisenberg-B-cycle-recall}
K^{(1)}(z+\tau|\tau) = K^{(1)}(z|\tau) - 2\pi i.
\end{equation}
On the torus, the fiberwise differential $\dfib$ involves
integrating $z_1$ over the fundamental domain.
The $B$-cycle quasi-periodicity means that the Green
function \emph{fails to be single-valued}: transporting
$z_1$ around the $B$-cycle shifts the propagator by the
constant $-2\pi i$. This failure is the source of the
curvature.

\emph{Step~3: Computing $\dfib^{\,2}$.}
Squaring the fiberwise differential on
$a^i \boxtimes a^j \cdot \eta_{12}^{\mathrm{ell}}$
amounts to: (i)~extract the OPE
residue (Step~1), then (ii)~contract the result with
a second propagator and take the $B$-cycle defect.
For the Heisenberg algebra, the only OPE contribution
is the double pole $\ell\delta^{ij}/(z-w)^2$.
The Laurent expansion of $K^{(1)}(z|\tau)$ near $z = 0$
is:
\[
K^{(1)}(z|\tau) = \frac{1}{z}
- \frac{\pi^2}{3}E_2(\tau)\,z + O(z^3).
\]
The local collision extraction is the same $d\log$ calculation as
in Proposition~\ref{prop:heisenberg-r-matrix}:
\[
\ell\delta^{ij}z^{-2}\,dz
=\ell\delta^{ij}z^{-1}\,d\log z .
\]
It produces the binary scalar residue
$r^{ij}(z)=\ell\delta^{ij}/z$. The genus-$1$ curvature is not a
new local residue. It is the global monodromy defect:
when $\dfib$ is squared, the $B$-cycle shift
$-2\pi i$ acts on the vacuum coefficient
$\ell\delta^{ij}\cdot|0\rangle$ from the first
application of $\dfib$.
Concretely, the Quillen anomaly formula gives:
\begin{equation}\label{eq:heisenberg-dfib-squared}
\dfib^{\,2}(a^i \boxtimes a^j
\cdot \eta_{12}^{\mathrm{ell}})
= \ell\delta^{ij}\cdot\omega_1
\end{equation}
where $\omega_1 = c_1(\mathbb{E})
\in H^2(\overline{\mathcal{M}}_{1,1})$
is the first Chern class of the Hodge bundle,
represented by the $B$-cycle monodromy class $(-2\pi i)$.

\emph{Step~4: Trace over generators.}
The genus-1 curvature of the \emph{full} rank-$d$
algebra is obtained by summing over the diagonal
self-contraction of each generator:
\begin{align}
m_0^{(ij)} &= \ell\delta^{ij}\cdot\omega_1,
\label{eq:rank-d-curvature}
\\
m_0 &= \sum_{i=1}^d m_0^{(ii)}
= d\ell\cdot\omega_1.
\end{align}
The obstruction coefficient is the scalar
prefactor:
\begin{equation}\label{eq:heisenberg-kappa-formula}
\boxed{\;\kappa(\mathcal{H}_\ell^{\oplus d})
= d\ell.\;}
\end{equation}

\emph{Step~5: Consistency checks.}
\begin{enumerate}[label=(\roman*)]
\item \emph{Central charge comparison.}
At level $\ell = 1$:
$\kappa(\mathcal{H}_1^{\oplus d}) = d = c$.
At general level the central charge remains $c=d$, while the
modular characteristic is $d\ell$.
\item \emph{Additivity.}
$\kappa(\mathcal{H}_\ell^{\oplus d_1}
\oplus \mathcal{H}_\ell^{\oplus d_2})
= (d_1 + d_2)\ell
= \kappa(\mathcal{H}_\ell^{\oplus d_1})
+ \kappa(\mathcal{H}_\ell^{\oplus d_2})$,
confirming the additivity of
Theorem~\ref{thm:modular-characteristic}(iii).
\item \emph{Koszul anti-symmetry.}
The Koszul dual $(\mathcal{H}_\ell)^!
= \operatorname{Sym}^{\mathrm{ch}}(V^*)$
(not $\mathcal{H}_{-\ell}$; see
Remark~\ref{rem:bosonization-not-koszul}); it has modular
characteristic $-\ell$ per generator, hence total characteristic
$-d\ell$. The free-field scalar pairing is
$d\ell+(-d\ell)=0$.
\item \emph{Gaussian shadow termination.}
Since the Heisenberg OPE has no simple pole,
the self-contraction loop producing the curvature
is the \emph{only} contribution to the shadow obstruction tower.
All higher shadows vanish ($\mathfrak{C} = 0$,
$\mathfrak{Q} = 0$), confirming shadow depth
$r_{\max} = 2$ (class~G).
\end{enumerate}

The first-principles residue and monodromy calculation gives
$\kappa(\mathcal{H}_\ell^{\oplus d}) = d\ell$
(Theorem~\ref{thm:heisenberg-obs}).

\emph{Comparison table.}
\begin{center}
\renewcommand{\arraystretch}{1.3}
\begin{tabular}{c|ccccccc}
$d$ & $c$ & $\kappa(\mathcal H_1^{\oplus d})$ & $F_1$ & $F_2$ & $F_3$ & $F_4$ & $F_5$ \\
\hline
$1$ & $1$ & $1$ & $1/24$ & $7/5760$ & $31/967680$ & $127/154828800$ & $73/3503554560$ \\
$2$ & $2$ & $2$ & $1/12$ & $7/2880$ & $31/483840$ & $127/77414400$ & $73/1751777280$ \\
$8$ & $8$ & $8$ & $1/3$ & $7/720$ & $31/120960$ & $127/19353600$ & $73/437944320$ \\
$16$ & $16$ & $16$ & $2/3$ & $7/360$ & $31/60480$ & $127/9676800$ & $73/218972160$ \\
$24$ & $24$ & $24$ & $1$ & $7/240$ & $31/40320$ & $127/6451200$ & $73/145981440$ \\
$26$ & $26$ & $26$ & $13/12$ & $91/2880$ & $403/483840$ & $1651/77414400$ & $949/1751777280$
\end{tabular}
\end{center}

At level~$1$, the case $d = 24$ gives $F_1 = 1$. This is the
Leech-lattice normalization of the rank-$24$ free-boson theory.

At level~$1$, the case $d = 26$ gives
$\kappa(\mathcal H_1^{\oplus 26}) = 26$, and the Virasoro subalgebra
has $\kappa_{\mathrm{Vir}} = c/2 = 13$
(Remark~\ref{rem:vir-bosonic-string-genus}). The factor-of-2
discrepancy: $26 = 2 \times 13$ reflects the difference between the
full Heisenberg algebra (spin-1 generators) and its Virasoro
subalgebra (spin-2 generator $T = \frac{1}{2}\normord{JJ}$).
\end{computation}


\subsection{\texorpdfstring{Eisenstein hierarchy for the rank-$d$ bar complex}{Eisenstein hierarchy for the rank-d bar complex}}

\begin{proposition}[Multi-boson elliptic coefficients;
\ClaimStatusProvedHere]
\label{prop:multi-boson-eisenstein}
\index{Eisenstein series!multi-boson corrections}
For the rank-$d$ Heisenberg algebra at common level~$\ell$, the
elliptic bar kernel is diagonal in the generator index:
\[
\langle a^i(z_1)a^j(z_2)\rangle_{E_\tau}^{\mathrm{bar}}
=\ell\delta^{ij}G_\tau(z_1-z_2),
\qquad
G_\tau(z)=\frac1z-\sum_{m\ge1}G_{2m}(\tau)z^{2m-1}.
\]
\begin{enumerate}[label=\textup{(\roman*)}]
\item the scalar genus-$1$ obstruction coefficient is
$d\ell$;
\item the first elliptic correction is
\begin{equation}\label{eq:multi-boson-E2}
d\ell\,G_2(\tau)
=d\ell\,\frac{\pi^2}{3}E_2(\tau);
\end{equation}
\item the higher local elliptic coefficients are
\begin{equation}\label{eq:multi-boson-E4}
d\ell\,G_{2m}(\tau)
=d\ell\,2\zeta(2m)E_{2m}(\tau)\qquad (m\ge2)
\end{equation}
before Wick multiplication by additional contraction edges.
\end{enumerate}
No universal scalar term $d\ell^2E_4(\tau)$ exists without first
choosing a graph, a bar degree, and a contraction pattern; powers of
$\ell$ count Wick contraction edges, while the trace over generator
indices contributes the factor~$d$ for each diagonal loop.
\end{proposition}

\begin{proof}
The diagonal OPE
$a^i(z)a^j(w)\sim\ell\delta^{ij}(z-w)^{-2}$ gives the bar residue
$\ell\delta^{ij}/z$ after $d\log$ extraction. Tensor factorization
then reduces the rank-$d$ kernel to the rank-one kernel times
$\delta^{ij}$. Taking the scalar trace over the diagonal generator
index gives $\sum_i\ell=d\ell$. The Laurent expansion of
$G_\tau$ gives $G_2=\pi^2E_2/3$ and
$G_{2m}=2\zeta(2m)E_{2m}$ for $m\ge2$, which proves the displayed
coefficients. Wick products multiply one copy of~$\ell$ per
contraction edge; there is no graph-independent scalar
$E_4$ correction at bar degree~$3$.
\end{proof}


\subsection{Genus-2 period-matrix carrier and scalar projection}

\begin{computation}[Heisenberg at genus 2: period matrix dependence]
\label{comp:heisenberg-genus2-period}
\index{Heisenberg algebra!genus 2}

At genus~$2$, the bar complex uses the propagator
$\eta_{ij}^{(2)} = d_z \log E(z_i, z_j|\Omega)$ where
$E(z,w|\Omega)$ is the prime form on the genus-2 Riemann surface
$\Sigma_2$ with period matrix $\Omega \in \mathfrak{H}_2$ (the
Siegel upper half-space of degree~2).

\emph{Prime form expansion.}
Near $z = w$, the prime form behaves as:
\[
E(z,w|\Omega) = (z-w)\bigl(1 + c_2(\Omega)(z-w)^2
+ c_4(\Omega)(z-w)^4 + \cdots\bigr)
\]
in a chosen local coordinate. The coefficients are local prime-form
jet tensors. After scalar contraction they map to the genus-$2$
Siegel modular-form ring; before that contraction they are not
canonical scalar modular forms.

\emph{Degree-2 curvature.}
The genus-2 scalar curvature is not the square of the local
one-form $\eta_{12}^{(2)}$ on the diagonal. It is the global
Hodge-class obstruction
\[
m_0^{(2)}=\kappa\,\lambda_2,
\qquad
\lambda_2=c_2(\mathbb E)\in H^4(\overline{\mathcal M}_2),
\]
obtained from the two independent $B$-cycle monodromy defects.
The free energy at genus~2 is:
\[
F_2(\mathcal{H}_\kappa) = \kappa \cdot \int_{\overline{\mathcal{M}}_{2,1}}
\psi_1^2 \lambda_2 = \kappa \cdot \frac{7}{5760}
\]
(Theorem~\ref{thm:heisenberg-all-genus}).

\emph{Period matrix integration.}
The period matrix
$\Omega = \bigl(\begin{smallmatrix} \tau_1 & \tau_{12} \\
\tau_{12} & \tau_2 \end{smallmatrix}\bigr)$ parameterizes the
complex structure of $\Sigma_2$ and supplies analytic representatives
for the prime-form kernel. The scalar bar free energy is the
tautological projection
\[
\int_{\overline{\mathcal{M}}_{2,1}}\psi_1^2\lambda_2=7/5760,
\]
not the integral of a pointwise coefficient $c_2(\Omega)$.

\emph{Explicit numerical values for rank-$d$ Heisenberg at genus~2.}
\begin{center}
\renewcommand{\arraystretch}{1.2}
\begin{tabular}{c|ccc}
$d$ & $F_2$ (exact) & $F_2$ (decimal) & Algebraic carrier \\
\hline
$1$ & $7/5760$ & $1.215 \times 10^{-3}$ & rank-$1$ Heisenberg \\
$2$ & $7/2880$ & $2.431 \times 10^{-3}$ & rank-$2$ Heisenberg \\
$8$ & $7/720$ & $9.722 \times 10^{-3}$ & $E_8$ lattice Heisenberg \\
$24$ & $7/240$ & $2.917 \times 10^{-2}$ & Leech lattice Heisenberg \\
$26$ & $91/2880$ & $3.160 \times 10^{-2}$ & rank-$26$ Heisenberg
\end{tabular}
\end{center}
\end{computation}


\subsection{Scalar bar free energy and analytic partition functions}

\begin{remark}[Systematic comparison at genus~\texorpdfstring{$g$}{g}]
\label{rem:bar-vs-partition-systematic}
\index{bar complex!vs.\ partition function}

The bar-complex free energy $F_g$ and the analytic partition
function $Z_g$ encode different information.
For the rank-$d$ Heisenberg at level $\kappa = 1$:

\begin{center}
\renewcommand{\arraystretch}{1.3}
\begin{tabular}{l|cc}
& Bar free energy $F_g$ & Partition function $Z_g$ \\
\hline
Type & Rational number & Function on $\mathcal{M}_g$ \\
Depends on & $d$ (rank) only & Full moduli $(\Sigma_g, \Omega)$ \\
$g = 1$ & $d/24$ & $(\operatorname{Im}\tau)^{-d/2}
|\eta(\tau)|^{-2d}$ \\
$g = 2$ & $7d/5760$ &
$[\det(\operatorname{Im}\Omega)]^{-d/2}
\cdot [F_2^{\mathrm{osc}}(\Omega)]^{-d}$ \\
Large $g$ & $\sim 2d/(2\pi)^{2g}$ &
$\sim [\det'(\Delta_g)]^{-d/2}$ \\
Growth & Bernoulli decay & $(2g)!$ growth \\
Convergence & Yes ($|x| < 2\pi$) & No (asymptotic)
\end{tabular}
\end{center}

The bar-complex free energy $F_g$ is a \emph{tautological intersection
number} on $\overline{\mathcal{M}}_g$, extracting the class
$\kappa \lambda_g \in H^{2g}(\overline{\mathcal{M}}_g)$
(Proposition~\ref{prop:bar-tautological-filtration}).
The partition function $Z_g$ is a function on $\mathcal{M}_g$,
encoding the full spectral data of the Laplacian on $\Sigma_g$.
They are not equal as mathematical objects. The scalar projection is
\[
F_g = \kappa \cdot \int_{\overline{\mathcal{M}}_{g,1}}
\psi_1^{2g-2} \lambda_g
\]
and the determinant-line partition function retains all non-tautological
spectral dependence. The comparison map is the projection from the
Quillen anomaly class to its $\lambda_g$ coefficient
(Remark~\ref{rem:verlinde-asymptotics}), not a scalar equality
between $F_g$ and $\log Z_g$.
\end{remark}


%%% ================================================================
%%% The Heisenberg shadow: exact Gaussianity
%%% ================================================================
\section{The Heisenberg shadow: exact Gaussianity}
\label{sec:heisenberg-shadow-gaussianity}

The genus expansion of \S\ref{sec:heisenberg-genus-expansion-eisenstein}
shows that the Heisenberg genus tower is controlled entirely by
$\kappa(\mathcal{H})$ and Bernoulli numbers. This section explains
\emph{why} from the perspective of the nonlinear shadow calculus
(Appendix~\ref{sec:nms-shadow-calculus}): the Heisenberg is the
Gaussian archetype; its complementarity potential is exactly
quadratic.

\subsection{Tower construction for the Heisenberg}
\label{subsec:heisenberg-tower-construction}
\index{shadow obstruction tower!Heisenberg construction}

Apply the universal tower construction
(Construction~\ref{constr:tower-template}) for
$\mathcal{H}_\kappa$.

\paragraph{Step~1: Degree~$2$.}
The double-pole OPE $J^i(z)\,J^j(w) \sim \kappa\,\delta^{ij}/(z-w)^2$
yields the Hessian $H_{\mathcal{H}} = \kappa\,x^2$ on
the cyclic slice. This gives
$\Theta_{\mathcal{H}}^{\leq 2} = \kappa(\mathcal{H})$.

\paragraph{Step~2: Degree~$3$, obstruction vanishes.}
The OPE has no simple pole ($\ell_2^{\mathrm{tr}} = 0$ on
$V_{\mathcal{H}}$), so the cubic shadow
$\mathfrak{C}_{\mathcal{H}} = 0$ identically.
The obstruction $o_3(\mathcal{H}) = 0$.
Therefore $\Theta_{\mathcal{H}}^{\leq 3}
= \Theta_{\mathcal{H}}^{\leq 2} = \kappa\,x^2$.

\paragraph{Step~3: Degree~$4$, obstruction vanishes by consequence.}
Since $\mathfrak{C} = 0$, the quartic obstruction
$o_4 = \tfrac12\{\mathfrak{C},\mathfrak{C}\}_H = 0$,
and the master equation forces $\mathfrak{Q}_{\mathcal{H}} = 0$.
Therefore $\Theta_{\mathcal{H}}^{\leq 4} = \kappa\,x^2$.

\paragraph{Step~4: Termination at degree~$2$.}
All higher obstructions vanish by induction
(no higher operations exist on the Heisenberg slice).
The tower terminates:
$\Theta_{\mathcal{H}}^{\leq r} = \kappa\,x^2$ for all $r \geq 2$.
This is exact Gaussianity.

\paragraph{Step~5: Arithmetic face.}
The connected sewing amplitudes are $a_\cH(N) = \sigma_{-1}(N) =
\sum_{d \mid N}d^{-1}$ (the sum-of-reciprocal-divisors function),
and the Dirichlet--sewing lift
(Definition~\ref{def:dirichlet-sewing-lift}) is
\begin{equation}\label{eq:heisenberg-euler-product-example}
S_\cH(u) = \zeta(u)\,\zeta(u{+}1)
= \prod_p \frac{1}{(1{-}p^{-u})(1{-}p^{-u-1})}\,.
\end{equation}
This is the archetypal exact Euler--Koszul case:
$W(\cH) = \{1\}$, so $\operatorname{ek}(\cH) = 0$
(Definition~\ref{def:euler-koszul-class}).
The sewing reproducing kernel
$K_\cH(s,t) = \zeta(s{+}t)\zeta(s{+}t{+}1)$
(Definition~\ref{def:sewing-kernel})
is a product of two zeta functions; its Gram matrix
$\mathscr{M}_\cH(\alpha) \succ 0$
by Theorem~\ref{thm:gram-positivity}.

The Heisenberg is the unique standard family
whose prime-side Li coefficients
(Definition~\ref{def:prime-side-li})
are positive at low order:
$\tilde\lambda_n(\cH) > 0$ for $1 \leq n \leq 6$,
$\tilde\lambda_7(\cH) < 0$
(Theorem~\ref{thm:li-asymptotics}). The eventual negativity
reflects the \emph{two-line spectral structure} of
$\Xi_\cH(u) = (u{-}1)\zeta(u)\zeta(u{+}1)$: its zeros lie
on two critical lines
$\operatorname{Re}(u) = 1/2$ and
$\operatorname{Re}(u) = -1/2$
(under GRH), and the displaced zeros force exponentially
negative contributions at large~$n$
(Proposition~\ref{prop:li-criterion-failure}).

\medskip
The formal statements are
Theorem~\ref{thm:heisenberg-exact-linearity} and
Corollary~\ref{cor:heisenberg-postnikov-termination} below;
the detailed OPE-level computations are in
\S\ref{subsec:heisenberg-shadow-jets}--\S\ref{subsec:heisenberg-exact-linearity}.

\subsection{Shadow jets of the Heisenberg}
\label{subsec:heisenberg-shadow-jets}

\begin{definition}[Shadow jets on a cyclic slice]
\label{def:heisenberg-shadow-jets}
\index{shadow jets!Heisenberg}
Let $V_{\mathcal{H}} \subset H^*(\mathrm{Def}_{\mathrm{cyc}}
(\mathcal{H}_\kappa), \ell_1)$ be the cyclic slice on which the
cyclic pairing restricts to a nondegenerate quadratic form
(Setup~\ref{setup:nms-cyclic-slice}). Write
$H_{\mathcal{H}} \colon V_{\mathcal{H}} \xrightarrow{\sim}
V_{\mathcal{H}}^*[-1]$ for the induced Hessian and
$P_{\mathcal{H}} := H_{\mathcal{H}}^{-1}$ for the propagator.
The \emph{quadratic}, \emph{cubic}, and \emph{quartic shadow jets}
of $\mathcal{H}_\kappa$ are the Taylor coefficients of the
complementarity potential
(Definition~\ref{def:nms-complementarity-potential}) on $V_{\mathcal{H}}$:
\begin{equation}\label{eq:heisenberg-shadow-expansion}
S_{\mathcal{H}}^{\le 4}(x)
\;=\;
\tfrac{1}{2}\,H_{\mathcal{H}}(x,x)
\;+\;
\tfrac{1}{3!}\,\mathfrak{C}_{\mathcal{H}}(x,x,x)
\;+\;
\tfrac{1}{4!}\,\mathfrak{Q}_{\mathcal{H}}(x,x,x,x).
\end{equation}
These are the degree-$2$, $3$, and $4$ projections of the universal
Maurer--Cartan class $\Theta_{\mathcal{H}}$
(Definition~\ref{def:nms-shadow-jets}). The quartic obstruction is
\begin{equation}\label{eq:heisenberg-quartic-obstruction}
\mathfrak{o}_{\mathcal{H}}^{(4)}
\;:=\;
\tfrac{1}{2}\,\{\mathfrak{C}_{\mathcal{H}},
\mathfrak{C}_{\mathcal{H}}\}_{H_{\mathcal{H}}}.
\end{equation}
The shadow jets satisfy the master equations
(Theorem~\ref{thm:nms-shadow-master-equations}):
\begin{equation}\label{eq:heisenberg-master-equations}
\nabla_{H_{\mathcal{H}}} \mathfrak{C}_{\mathcal{H}} = 0,
\qquad
\nabla_{H_{\mathcal{H}}} \mathfrak{Q}_{\mathcal{H}}
+ \mathfrak{o}_{\mathcal{H}}^{(4)} = 0,
\end{equation}
where $\nabla_{H_{\mathcal{H}}}(X) :=
\{H_{\mathcal{H}}, X\}_{H_{\mathcal{H}}}$.
\end{definition}

\subsection{Exact linearity and vanishing of the nonlinear jets}
\label{subsec:heisenberg-exact-linearity}

\begin{theorem}[Heisenberg exact linearity; \ClaimStatusProvedHere]
\label{thm:heisenberg-exact-linearity}
\index{Heisenberg algebra!exact linearity}
On the Heisenberg deformation line one has
\begin{equation}\label{eq:heisenberg-shadow-vanishing}
\mathfrak{C}_{\mathcal{H}} = 0,
\qquad
\mathfrak{o}_{\mathcal{H}}^{(4)} = 0,
\qquad
\mathfrak{Q}_{\mathcal{H}} = 0.
\end{equation}
Equivalently, through quartic order the universal class reduces to its
quadratic part:
\begin{equation}\label{eq:heisenberg-theta-quadratic}
\Theta_{\mathcal{H}}^{\le 4} = H_{\mathcal{H}}.
\end{equation}
\end{theorem}

\begin{proof}
The Heisenberg deformation problem is governed by a constant Poisson
structure: the OPE $J^i(z)\,J^j(w) \sim \kappa\,\delta^{ij}/(z-w)^2$
has no simple-pole term, so the corresponding $L_\infty$ structure on
the cyclic deformation complex carries $\ell_1^{\mathrm{tr}} = 0$
(the transferred differential vanishes on $V_{\mathcal{H}}$) and no
higher operations $\ell_n^{\mathrm{tr}}$ for $n \ge 2$ beyond the
pairing $H_{\mathcal{H}}$ itself. On the bar side this is reflected
in the absence of operations $m_n$ with $n \ge 3$ on the
Heisenberg slice. Therefore the cubic shadow
$\mathfrak{C}_{\mathcal{H}}(x,y,z) =
\sum_{\mathrm{cyc}} \langle x, \ell_2^{\mathrm{tr}}(y,z)\rangle$
vanishes identically, the quartic obstruction
$\mathfrak{o}_{\mathcal{H}}^{(4)} =
\frac{1}{2}\{\mathfrak{C}_{\mathcal{H}},
\mathfrak{C}_{\mathcal{H}}\}_{H_{\mathcal{H}}}$
vanishes as a consequence, and the master
equation~\eqref{eq:heisenberg-master-equations} then forces
$\mathfrak{Q}_{\mathcal{H}} = 0$.
\end{proof}

\begin{corollary}[Heisenberg shadow obstruction tower:
finite termination at degree~$2$; \ClaimStatusProvedHere]
\label{cor:heisenberg-postnikov-termination}
\index{shadow obstruction tower!Heisenberg termination}
The shadow obstruction tower
\textup{(}Definition~\textup{\ref{def:shadow-postnikov-tower})}
of the Heisenberg algebra terminates at degree~$2$:
$\Theta_{\mathcal{H}}^{\leq r} = \Theta_{\mathcal{H}}^{\leq 2}
= \kappa(\mathcal{H})$ for all $r \geq 2$,
with $o_{r+1}(\mathcal{H}) = 0$ for all $r \geq 2$.
The Heisenberg tower is exactly Gaussian.
(The Heisenberg spectral $R$-matrix $R(z) = z^{2h}$ is the simplest case of the closed-form Virasoro $R$-matrix $R(z) = z^{2h}\exp(-(c/4)/z^2)$ of Computation~\ref{comp:virasoro-spectral-r-matrix}.)
\end{corollary}

\begin{proof}
Immediate from Theorem~\ref{thm:heisenberg-exact-linearity}:
the vanishing of all higher shadows means
every obstruction class vanishes, so the tower stabilizes
at the scalar level.
\end{proof}

\begin{computation}[OPE derivation of Heisenberg shadow vanishing]
\label{comp:heisenberg-shadow-ope-derivation}
The explicit OPE-level computations underlying
Theorem~\ref{thm:heisenberg-exact-linearity} are as follows.
Compare with the nonvanishing cubic and quartic shadows of the affine and
Virasoro families (Chapter~\ref{chap:kac-moody},
Chapter~\ref{chap:w-algebras}).

\medskip
\noindent\textbf{Step 1: The cyclic slice.}
For the rank-$1$ Heisenberg $\mathcal{H}_\kappa$ with current $J$
and OPE
\begin{equation}\label{eq:heisenberg-ope-shadow-section}
 J(z)\,J(w) \;\sim\; \frac{\kappa}{(z-w)^2},
\end{equation}
the cyclic deformation complex
$(\mathrm{Def}_{\mathrm{cyc}}(\mathcal{H}_\kappa), \ell_1)$ has a
one-dimensional cohomological slice
\begin{equation}\label{eq:heisenberg-cyclic-slice}
 V_{\mathcal{H}} \;=\; \mathbb{C}\langle J \rangle,
\end{equation}
where $J$ represents the marginal deformation direction (level
shift $\kappa \mapsto \kappa + \epsilon$). The cyclic pairing on
mode space is
\begin{equation}\label{eq:heisenberg-cyclic-pairing}
 \langle J_m,\, J_n \rangle_{\mathrm{cyc}}
 \;=\; \kappa\,\delta_{m,-n},
\end{equation}
extracted from the leading-pole coefficient: $\operatorname{Res}_{z=w}
(z-w)\,\langle J(z)\,J(w)\rangle = \kappa$.

\medskip
\noindent\textbf{Step 2: The Hessian from the OPE.}
The Hessian $H_{\mathcal{H}} \colon
V_{\mathcal{H}} \to V_{\mathcal{H}}^*[-1]$ is the restriction of
the cyclic form to $V_{\mathcal{H}}$.
Writing $x \in V_{\mathcal{H}}$ for the coordinate dual to $J$:
\begin{equation}\label{eq:heisenberg-hessian-explicit}
 H_{\mathcal{H}}(x,x) \;=\; \kappa\,x^2.
\end{equation}
The propagator is its inverse:
\begin{equation}\label{eq:heisenberg-propagator-explicit}
 P_{\mathcal{H}} \;=\; H_{\mathcal{H}}^{-1}
 \;=\; \kappa^{-1}.
\end{equation}

\medskip
\noindent\textbf{Step 3: Vanishing of the cubic shadow, from the
 OPE.}
The cubic shadow is
\begin{equation}\label{eq:heisenberg-cubic-formula}
 \mathfrak{C}_{\mathcal{H}}(x,y,z)
 \;=\; \sum_{\mathrm{cyc}} \bigl\langle x,\,
 \ell_2^{\mathrm{tr}}(y,z) \bigr\rangle_{\mathrm{cyc}},
\end{equation}
where $\ell_2^{\mathrm{tr}}$ is the transferred binary operation on
$V_{\mathcal{H}}$. This operation is read off the simple-pole
coefficient of the OPE: writing
\begin{equation}\label{eq:ope-simple-pole-extraction}
 J(z)\,J(w)
 \;=\;
 \frac{\kappa}{(z-w)^2}
 \;+\;
 \frac{[J,J](w)}{z-w}
 \;+\; \text{regular},
\end{equation}
the $(z-w)^{-1}$ residue is the current-algebra Lie bracket
$[J,J]$ restricted to the cyclic slice. For the Heisenberg, this
coefficient is \emph{identically zero}: the OPE has no simple
pole. Therefore:
\begin{align}\label{eq:heisenberg-ell2-vanishing}
 \ell_2^{\mathrm{tr}}(J,J) &= 0
 \quad\text{on } V_{\mathcal{H}},\\
 \label{eq:heisenberg-cubic-vanishing}
 \mathfrak{C}_{\mathcal{H}} &= 0.
\end{align}
This is the decisive step: the vanishing is not a
cancellation; it is the \emph{absence} of a Lie bracket in the
current algebra.

\medskip
\noindent\textbf{Step 4: Vanishing of the quartic shadow, from
 homotopy transfer.}
The quartic shadow is
\begin{equation}\label{eq:heisenberg-quartic-formula}
 \mathfrak{Q}_{\mathcal{H}}(x_1, x_2, x_3, x_4)
 \;=\;
 \sum_{\mathrm{cyc}} \bigl\langle x_1,\,
 \ell_3^{\mathrm{tr}}(x_2, x_3, x_4)
 \bigr\rangle_{\mathrm{cyc}},
\end{equation}
where the transferred ternary operation $\ell_3^{\mathrm{tr}}$
arises from the $A_\infty$ homotopy transfer formula
(Appendix~\ref{app:homotopy-transfer}):
\begin{equation}\label{eq:heisenberg-transfer-formula}
 \ell_3^{\mathrm{tr}}
 \;=\;
 \ell_2^{\mathrm{tr}} \circ (h \otimes \mathrm{id})
 \circ \ell_2^{\mathrm{tr}}
 \;+\;
 \ell_2^{\mathrm{tr}} \circ (\mathrm{id} \otimes h)
 \circ \ell_2^{\mathrm{tr}}
 \;+\; (\text{terms involving } \ell_1^{\mathrm{tr}}),
\end{equation}
where $h$ is the contracting homotopy on the acyclic complement.
Since $\ell_2^{\mathrm{tr}} = 0$ on $V_{\mathcal{H}}$
(equation~\eqref{eq:heisenberg-ell2-vanishing}), every composite in
the transfer formula has at least one factor of
$\ell_2^{\mathrm{tr}}|_{V_{\mathcal{H}}}$ and therefore vanishes:
\begin{align}\label{eq:heisenberg-ell3-vanishing}
 \ell_3^{\mathrm{tr}}(J,J,J) &= 0
 \quad\text{on } V_{\mathcal{H}},\\
 \label{eq:heisenberg-quartic-vanishing}
 \mathfrak{Q}_{\mathcal{H}} &= 0.
\end{align}
The same argument applies at all degrees: the transferred operations
$\ell_n^{\mathrm{tr}}$ vanish on $V_{\mathcal{H}}$ for all
$n \ge 2$, since the transfer tree formulae are always built from
composites of lower-degree operations.

\medskip
\noindent\textbf{Step 5: The complementarity potential.}
The full complementarity potential
(Definition~\ref{def:nms-complementarity-potential}) is therefore
exactly quadratic:
\begin{equation}\label{eq:heisenberg-complementarity-potential}
 S_{\mathcal{H}}(x)
 \;=\;
 \tfrac{1}{2}\,\kappa\,x^2.
\end{equation}
All higher jets $\mathfrak{C}_{\mathcal{H}},
\mathfrak{Q}_{\mathcal{H}}, \mathfrak{o}^{(5)}_{\mathcal{H}},
\ldots$ vanish. The dual Lagrangian is
\begin{equation}\label{eq:heisenberg-dual-lagrangian}
 M_{\mathcal{H}^!}
 \;=\;
 \operatorname{graph}(dS_{\mathcal{H}})
 \;=\;
 \{(x, \kappa x)\} \;\subset\;
 V_{\mathcal{H}} \oplus V_{\mathcal{H}}^*,
\end{equation}
which is a \emph{linear} Lagrangian submanifold. Hence the
Heisenberg satisfies the fake complementarity criterion
(Proposition~\ref{prop:fake-complementarity-criterion}): the
complementarity is ``fake'' in the precise sense that the dual
Lagrangian is the graph of a \emph{linear} map. The Heisenberg is
the archetype of fake complementarity.

\medskip
\noindent\textbf{Step 6: Bridge to the Bernoulli genus tower.}
When $S_{\mathcal{H}}$ is exactly quadratic, the genus recursion
$\xi^*\Theta = \Theta \star \Theta$
(Theorem~\ref{thm:nms-separating-boundary-recursion}) involves only
the Hessian: every boundary term contracts to
\begin{equation}\label{eq:heisenberg-hessian-contraction}
 H_{\mathcal{H}} \star H_{\mathcal{H}}
 \;=\;
 \operatorname{tr}(P_{\mathcal{H}} \circ H_{\mathcal{H}})
 \;=\;
 \kappa \cdot \kappa^{-1}
 \;=\; 1,
\end{equation}
which is the identity under propagator contraction on a rank-$1$
space. The genus-$g$ free energy is therefore purely tautological :
\begin{equation}\label{eq:heisenberg-genus-g-tautological}
 F_g(\mathcal{H}_\kappa)
 \;=\;
 \kappa \cdot \lambda_g^{\mathrm{FP}},
 \qquad
 \lambda_g^{\mathrm{FP}}
 \;=\;
 \frac{2^{2g-1}-1}{2^{2g-1}} \cdot \frac{|B_{2g}|}{(2g)!}.
\end{equation}
where $\lambda_g^{\mathrm{FP}}$ is the Faber--Pandharipande
tautological constant and $B_{2g}$ is the $(2g)$-th Bernoulli
number. The first values are:
\begin{align}
 \label{eq:heisenberg-F1-shadow}
 F_1 &= \kappa \cdot \frac{1}{2} \cdot \frac{|B_2|}{2!}
 \;=\; \kappa \cdot \frac{1}{2} \cdot \frac{1/6}{2}
 \;=\; \frac{\kappa}{24},\\[4pt]
 \label{eq:heisenberg-F2-shadow}
 F_2 &= \kappa \cdot \frac{7}{8} \cdot \frac{|B_4|}{4!}
 \;=\; \kappa \cdot \frac{7}{8} \cdot \frac{1/30}{24}
 \;=\; \frac{7\kappa}{5760},\\[4pt]
 \label{eq:heisenberg-F3-shadow}
 F_3 &= \kappa \cdot \frac{31}{32} \cdot \frac{|B_6|}{6!}
 \;=\; \kappa \cdot \frac{31}{32} \cdot \frac{1/42}{720}
 \;=\; \frac{31\kappa}{967680}.
\end{align}
The generating function assembles into the $\hat{A}$-genus:
\begin{equation}\label{eq:heisenberg-ahat-from-shadows}
 \sum_{g \ge 1} F_g(\mathcal{H}_\kappa)\,\hbar^{2g-2}
 \;=\;
 \frac{\kappa}{\hbar^2}
 \bigl(\hat{A}(i\hbar) - 1\bigr),
\end{equation}
confirming Theorem~\ref{thm:family-index} directly from the shadow
structure. The exactness of the quadratic potential is the shadow-
calculus explanation for the $\hat{A}$-genus: the Heisenberg genus
tower is Gaussian because its complementarity potential \emph{is}
Gaussian.
\end{computation}

\begin{corollary}[Gaussian boundary law; \ClaimStatusConditional]
\label{cor:heisenberg-gaussian-boundary}
\index{Heisenberg algebra!Gaussian boundary law}
For the Heisenberg family the separating boundary recursion is purely
Gaussian:
\begin{equation}\label{eq:heisenberg-gaussian-boundary}
\xi^*H_{\mathcal{H}}
= H_{\mathcal{H}} \star H_{\mathcal{H}},
\qquad
\xi^*\mathfrak{C}_{\mathcal{H}} = 0,
\qquad
\xi^*\mathfrak{Q}_{\mathcal{H}} = 0,
\end{equation}
where $\xi$ is any separating clutching map and $\star$ is the
single-edge sewing product
(Construction~\ref{constr:nms-sewing-product}).
\end{corollary}

\begin{proof}
The separating sewing formulae
(Theorem~\ref{thm:nms-separating-boundary-recursion}) express
$\xi^*\mathfrak{C}$ and $\xi^*\mathfrak{Q}$ as polynomial
expressions in $\mathfrak{C}$, $\mathfrak{Q}$,
$\mathfrak{o}^{(4)}$, and $H$ via the propagator $P$. Since
$\mathfrak{C}_{\mathcal{H}} = 0$ and
$\mathfrak{Q}_{\mathcal{H}} = 0$
(Theorem~\ref{thm:heisenberg-exact-linearity}),
all nonlinear terms in the recursion vanish. Only the quadratic
sewing $\xi^*H_{\mathcal{H}} = H_{\mathcal{H}} \star H_{\mathcal{H}}$
survives.
\end{proof}

\subsection{The Gaussian archetype}
\label{subsec:heisenberg-gaussian-archetype}

\begin{remark}[Exact Gaussianity and the scalar genus tower]
\label{rem:heisenberg-exact-gaussianity}
\index{Heisenberg algebra!Gaussianity}
The Heisenberg is the unique frame family whose complementarity
potential is exactly quadratic through all orders. This is the
algebraic reason the Heisenberg genus tower
is determined by $\kappa$ alone: the scalar modular shadow
characteristic
\[
\mathrm{ch}_{\mathrm{mod}}(\mathcal{H}_\kappa;\hbar)
\;=\;
\sum_{g \ge 0} \hbar^g\,\kappa(\mathcal{H}_\kappa)\,\lambda_g
\]
collapses to its scalar trace. The higher-degree jets that would
carry non-scalar modular information (the cubic $\mathfrak{C}$
governing tree-level boundary contributions, the quartic
$\mathfrak{Q}$ encoding contact corrections) all vanish.
There is no nonlinear modular data beyond the $\hat{A}$-genus.
\end{remark}

\begin{remark}[Phase index]
\label{rem:heisenberg-phase-index}
\index{nonlinear phase index!Heisenberg}
The nonlinear phase index
(Definition~\ref{def:nms-nonlinear-phase-index}) takes the value
$\nu_{\mathrm{nl}}(\mathcal{H}) = \infty$ for the Heisenberg: the
complementarity potential has no nonlinear jet at any order. This is
the Gaussian archetype of the trichotomy
(Theorem~\ref{thm:nms-archetype-trichotomy}). In the classification
of the five standard families
(Proposition~\ref{conj:nms-nonlinear-phase-standard}), the Heisenberg
is the only family with $\nu_{\mathrm{nl}} = \infty$; the affine and
Virasoro families have $\nu_{\mathrm{nl}} = 3$ (cubic leading
nonlinearity), while the $\beta\gamma$ system has
$\nu_{\mathrm{nl}} = 4$ (quartic contact). The Heisenberg shadow
tower terminates at degree~$2$
(Theorem~\ref{thm:nms-finite-termination}), producing the simplest
possible modular package.
\end{remark}

\begin{remark}[Completion kinematics of the Heisenberg]
\label{rem:heisenberg-completion-kinematics}
\index{completion kinematics!Heisenberg}
\index{completion entropy!Heisenberg}
The Heisenberg is kinematically trivial: $G_{\mathcal{H}_k}(t) = t$
(one primitive of reduced weight~$1$),
$H_{\mathcal{H}_k}(t) = t/(1-t)$, $\rho_K = 1$, $h_K = 0$,
$\Delta_{\mathcal{H}_k}(t) = 0$. The reduced-weight-$q$ window
$K_q(\mathcal{H}_k)$ is one-dimensional for every~$q$ (a single
bar word $\partial^{q-1}\alpha$), and the bar differential on each
window is zero. This is the unique G-class family
(Table~\ref{tab:completion-kinematics}): both the shadow depth
(degree~$2$) and the completion entropy ($h_K = 0$) are minimal.
The kinematic triviality of the Heisenberg is the algebraic reason
that all other families can be calibrated against it: every
completion invariant is measured relative to this zero point.

In the holographic modular Koszul datum
$\mathcal{H}(T) = (\cA, \cA^!, C, r^{\mathrm{coll}}(z), \Theta_\cA, \nabla^{\mathrm{hol}})$
of the concordance
(\S\ref{subsec:completion-kinematics-theory}),
the Heisenberg specialisation is: $\cA = \mathcal{H}_k$,
$\cA^! = \operatorname{Sym}^{\mathrm{ch}}(V^*)$ with $\kappa(\cA^!) = -k$ (the chiral symmetric algebra, not $\cH_{-k}$),
$C = $ direct-sum complementarity,
$r^{\mathrm{Heis}}_k(z) = k/z$ (the abelian $r$-matrix),
$\pi_{\mathrm{scal}}(\Theta_{\mathcal{H}_k})
= k \cdot \eta \otimes \Lambda$
(scalar projection, no higher shadow corrections),
$\nabla^{\mathrm{hol}} = d$ (flat, since the shadow vanishes).
Every subsequent family is measured against this datum.
\end{remark}

\subsection{The Heisenberg primitive kernel}
\label{subsec:heisenberg-primitive-kernel}
\index{primitive kernel!Heisenberg}

\begin{proposition}[Heisenberg primitive kernel;
\ClaimStatusConditional]
\label{prop:heisenberg-primitive-kernel}
The primitive logarithmic modular kernel
\textup{(}Definition~\textup{\ref{def:primitive-log-modular-kernel})}
of the Heisenberg algebra~$\mathcal{H}_\kappa$ is
\begin{equation}\label{eq:heisenberg-primitive-kernel}
\mathfrak{K}_{\mathcal{H}}
\;=\;
K_{0,2}^{\mathcal{H}} + K_{1,1}^{\mathcal{H}},
\end{equation}
where $K_{0,2}^{\mathcal{H}} = \kappa\,\delta^{ij}$ is the degree-$2$
corolla encoding the double-pole OPE and
$K_{1,1}^{\mathcal{H}} = \Delta_{\mathrm{ns}}(K_{0,2}^{\mathcal{H}})$
is the genus-$1$ corolla from non-separating self-sewing of
$K_{0,2}$. All higher corollas vanish:
\begin{equation}\label{eq:heisenberg-higher-corollas-vanish}
K_{0,n}^{\mathcal{H}} = 0 \;\; (n \geq 3),
\qquad
R_\rho^{\mathcal{H}} = 0 \;\;\text{for all rigid types } \rho.
\end{equation}
\end{proposition}

\begin{proof}
The absence of a simple pole in the Heisenberg OPE
$J^i(z)\,J^j(w) \sim \kappa\,\delta^{ij}/(z{-}w)^2$
gives $\ell_2^{\mathrm{tr}} = 0$ on the cyclic slice, so the
transferred ternary and all higher operations vanish:
$\ell_n^{\mathrm{tr}} = 0$ for $n \geq 2$. Therefore
$K_{0,n}^{\mathcal{H}} = 0$ for $n \geq 3$. The rigid-cutting
residues $R_\rho$ arise from planted-forest strata of the
logarithmic FM boundary
(Theorem~\ref{thm:logfm-modular-cocomposition}), which involve
compositions of transferred operations through at least one
ternary vertex; since all such operations vanish, $R_\rho = 0$.
\end{proof}

\begin{proposition}[Heisenberg primitive shell equations;
\ClaimStatusConditional]
\label{prop:heisenberg-primitive-shell}
\index{shell equation!Heisenberg}
The primitive master equation
$d\,\mathfrak{K}_{\mathcal{H}}
+ \mathfrak{K}_{\mathcal{H}} \star \mathfrak{K}_{\mathcal{H}}
+ \hbar\,\Delta_{\mathrm{ns}}(\mathfrak{K}_{\mathcal{H}}) = 0$
\textup{(}Theorem~\textup{\ref{thm:primitive-to-global-reconstruction})}
reduces on the Heisenberg to the single channel:
\begin{equation}\label{eq:heisenberg-shell-genus2}
R_g(\mathfrak{K}_{\mathcal{H}})
\;=\;
\Delta_{\mathrm{ns}}(K_{g-1,1}^{\mathcal{H}})
\qquad (g \geq 2).
\end{equation}
The separating bracket $[K_{g_1}, K_{g_2}] = 0$ \textup{(}rank~$1$,
antisymmetry\textup{)}, and the rigid-cutting channel is empty.
The genus tower is generated by iterated BV self-contraction alone.
\end{proposition}

\begin{proof}
The separating channel
$\frac{1}{2}[K_{g_1,1}, K_{g_2,1}]$ with $g_1 + g_2 = g$
vanishes for the Heisenberg: each $K_{g,1}$ is a scalar
(rank-$1$ branch space), and the pre-Lie bracket antisymmetrises
two copies of the same scalar, giving zero. The rigid-cutting
channel is empty since $R_\rho = 0$
(Proposition~\ref{prop:heisenberg-primitive-kernel}). Only the
BV self-contraction $\Delta_{\mathrm{ns}}$ survives.
\end{proof}

\begin{computation}[Heisenberg branch BV action]
\label{comp:heisenberg-branch-bv}
\index{branch BV action!Heisenberg}
The branch-active quotient
\textup{(}Definition~\textup{\ref{def:reduced-branch-master-action})}
is one-dimensional: $V_{\mathcal{H}}^{\mathrm{br}} = \mathbb{C}$
with coordinate~$x$. The branch master action is purely quadratic:
\begin{equation}\label{eq:heisenberg-branch-action}
S_{\mathcal{H}}^{\mathrm{br}}(x)
\;=\;
\tfrac{1}{2}\,\kappa\,x^2.
\end{equation}
The branch QME
\textup{(}Proposition~\textup{\ref{prop:branch-master-equation})}
is trivially satisfied: the master action is Gaussian, so
$\{S,S\}_{\mathrm{BV}} = 0$ in the BV bracket on
$\widehat{\mathrm{Sym}}((V^+)^\vee)$.
The spectral discriminant and metaplectic half-density
\textup{(}Corollary~\textup{\ref{cor:metaplectic-square-root})}
are:
\begin{equation}\label{eq:heisenberg-metaplectic}
\Delta_{\mathcal{H}}(x) = 1 - \kappa\,x,
\qquad
\delta_{\mathcal{H}}(x) = (1 - \kappa\,x)^{1/2}.
\end{equation}
Both are linear/square-root: the Heisenberg branch BV packet is a
free rank-$1$ BV algebra with Gaussian measure. The flat connection
on the primitive deformation complex
\textup{(}\S\textup{\ref{sec:primitive-kernel-intro})}
is $\nabla^{\mathrm{mod}}_{\mathcal{H}} = d$: trivially flat,
since $[\Theta_{\mathcal{H}}, -] = 0$ on all higher corollas.
\end{computation}

\section{Modular cooperad structure on the open sector}
\label{sec:heisenberg-modular-cooperad}
\index{modular cooperad!Heisenberg|textbf}
\index{Heisenberg!modular cooperad|textbf}

The sewing amplitudes of the preceding sections assemble genus by
genus, but a genus-$g$ amplitude alone does not know how it decomposes
under degeneration to a nodal curve. The global modular cooperad
completion
(\S\ref{sec:concordance-modular-cooperad-completion}) asks for a
system of cocomposition maps $\{\delta_\Gamma\}$ indexed by stable
bordered graphs, satisfying coassociativity, self-sewing = trace, and
compatibility with the closed-sector MC element
$\Theta^{\mathrm{oc}}$. For a general chiral algebra, verifying these
axioms requires the full derived structure of the open-sector
category. In the Heisenberg class-G Fock packet, the one-particle
Bergman reduction
(Theorem~\ref{thm:heisenberg-one-particle-sewing}) reduces the
restricted trace-class verification to operator algebra on the
Bergman space $A^2(D)$, where coassociativity is the multiplicativity
of Fredholm determinants.

\subsection{The open-sector category}
\label{subsec:heisenberg-open-sector-category}

\begin{proposition}[Completed Fock open sector for Heisenberg;
\ClaimStatusProvedHere]
\label{prop:heisenberg-open-sector}
\index{Heisenberg!open-sector category}
Let $\Cop(\cH_\kappa)$ denote the completed Fock open sector: the
idempotent-complete dg subcategory generated by finite direct sums of
Fock modules and completed in conformal weight for trace-class
sewing. It has the following structure.
\begin{enumerate}[label=\textup{(\roman*)}]
\item \emph{Objects.} The Fock modules
 $\cF_\lambda$ ($\lambda \in \mathbb{C}$), generated by a
 highest-weight vector $|\lambda\rangle$ with
 $\alpha_0\,|\lambda\rangle = \lambda\,|\lambda\rangle$
 and $\alpha_n\,|\lambda\rangle = 0$ for $n > 0$.
\item \emph{Semisimplicity in the Fock envelope.} Objects in this
 completed Fock sector are finite direct sums of the $\cF_\lambda$.
 No assertion is made here about arbitrary $\cH_\kappa$-modules.
\item \emph{Dualizability.} Each $\cF_\lambda$ is dualizable
 in the completed trace-class sense: the dual is
 $\cF_\lambda^\vee \simeq \cF_{-\lambda}$ under the contragredient
 pairing, and the evaluation and coevaluation maps are induced by the
 Shapovalov form. Weight spaces are finite-dimensional, and the
 $q$-completed trace converges for sewing parameters $|q|<1$.
\end{enumerate}
\end{proposition}

\begin{proof}
Part~(i): This is the standard classification of Fock representations of
the Heisenberg algebra~\cite{FBZ04}. The $\alpha_n$ ($n \in \mathbb{Z}$)
act on $\cF_\lambda = \mathbb{C}[x_1, x_2, \ldots] \otimes
|\lambda\rangle$ by $\alpha_{-n} = x_n$ (creation, $n > 0$),
$\alpha_n = n\kappa\,\partial_{x_n}$ (annihilation, $n > 0$), and
$\alpha_0 = \lambda$.

Part~(ii): The statement is by definition a statement about the Fock
envelope. The Heisenberg OPE has no simple-pole bracket, so the
transferred $A_\infty$-module operations that would mix Fock summands
vanish for $n\ge3$. This gives a semisimple completed envelope on the
chosen Fock generators; it is not a classification of all
$\cH_\kappa$-modules.

Part~(iii): Dualizability of the simple objects in a semisimple
category is automatic: the contragredient module $\cF_\lambda^\vee$
is generated by $\langle\lambda|$ with $\langle\lambda|\,\alpha_0
= \lambda\,\langle\lambda|$, $\langle\lambda|\,\alpha_{-n} = 0$
for $n > 0$. The Shapovalov form
$\langle\lambda|\,\alpha_{n_1} \cdots \alpha_{n_k}\,
\alpha_{-m_1} \cdots \alpha_{-m_l}\,|\lambda\rangle$ is
nondegenerate on each weight space, so the evaluation
$\operatorname{ev}\colon \cF_\lambda^\vee \otimes \cF_\lambda
\to \mathbb{C}$ and coevaluation
$\operatorname{coev}\colon \mathbb{C} \to \cF_\lambda \otimes
\cF_\lambda^\vee$ satisfy the snake identities on the completed Fock
sector. Finite-dimensionality of each conformal-weight space gives
trace-class $q$-traces for $|q|<1$; uncompleted Fock spaces are not
proper finite-dimensional objects.
\end{proof}

\subsection{The sewing cocomposition maps}
\label{subsec:heisenberg-sewing-cocomposition}

\begin{construction}[Cocomposition maps from sewing]
\label{constr:heisenberg-cocomposition}
\index{cocomposition!Heisenberg}
\index{stable bordered graph}
Let $\Gamma$ be a stable bordered graph of type $(g, n)$:
a graph with vertices $v \in V(\Gamma)$ decorated by genera
$g_v \geq 0$, legs $L(\Gamma)$ with $|L(\Gamma)| = n$, and
internal edges $E(\Gamma)$ encoding degenerations, subject to
the stability condition $2g_v - 2 + n_v > 0$ at each vertex
(with $n_v$ the valence of $v$, counting both half-edges and
legs) and the genus identity $g = \sum_v g_v + |E(\Gamma)| -
|V(\Gamma)| + 1$.

For each such $\Gamma$, define the cocomposition map
\begin{equation}\label{eq:heisenberg-cocomposition}
\delta_\Gamma
\;\colon\;
\operatorname{End}_{\Cop}(\cF_\lambda)
\;\longrightarrow\;
\bigotimes_{v \in V(\Gamma)}
\operatorname{End}_{\Cop}(\cF_\lambda)
\end{equation}
by the sewing formula: decompose the genus-$g$ amplitude into
a product of vertex amplitudes connected by sewing operators.
At the one-particle level
\textup{(}Theorem~\textup{\ref{thm:heisenberg-one-particle-sewing}}\textup{)},
each internal edge $e \in E(\Gamma)$ contributes a sewing
operator $\sewop_{q_e}$ on the Bergman space $A^2(D)$, and each
vertex $v$ contributes the second quantization
$\secquant(R_v)$ of its vertex restriction operator.

Concretely, for a genus-$g$ amplitude
$Z \in \operatorname{End}_{\Cop}(\cF_\lambda)$, the
cocomposition is:
\begin{equation}\label{eq:heisenberg-delta-gamma-formula}
\delta_\Gamma(Z)
\;=\;
\bigotimes_{v \in V(\Gamma)}
Z_v,
\qquad
Z_v
= \operatorname{Tr}_{\{e \,:\, e \text{ at } v\}}
\!\Bigl(
\secquant\!\Bigl(
\prod_{e \text{ at } v} \sewop_{q_e}
\Bigr)^{\!\kappa}\,\Bigr),
\end{equation}
where the trace is over the one-particle Hilbert spaces
associated to the internal half-edges at vertex~$v$, and the
product of sewing operators acts on the tensor product of
the corresponding Bergman spaces.

At the one-particle level, this becomes:
\begin{equation}\label{eq:heisenberg-one-particle-cocomposition}
\delta_\Gamma^{(1)}
\;\colon\;
\sewker_g
\;\longmapsto\;
\bigotimes_{v \in V(\Gamma)}
\sewker_{g_v, n_v},
\end{equation}
where $\sewker_{g_v, n_v}$ is the one-particle sewing kernel of the
surface $\Sigma_{g_v, n_v}$ at vertex $v$, and the internal sewing
parameters $q_e$ mediate the decomposition.
\end{construction}

\begin{theorem}[CT-$2$ for Heisenberg: modular cooperad on $\Cop(\cH_\kappa)$;
\ClaimStatusConditional]
\label{thm:heisenberg-modular-cooperad}
\label{thm:ct2-heisenberg}
\index{modular cooperad!Heisenberg!main theorem}
\index{CT-2!Heisenberg|textbf}
Assume the completed Fock open sector of
Proposition~\ref{prop:heisenberg-open-sector} and sewing parameters
$|q_e|<1$ on every internal edge. Then the completed Fock packet inside
$\Cop(\cH_\kappa)$ carries a trace-class modular cooperad structure
$\{\delta_\Gamma\}_{\Gamma \in \mathcal{G}^{\mathrm{bord}}}$
indexed by stable bordered graphs
\textup{(}Construction~\textup{\ref{constr:heisenberg-cocomposition}}\textup{)},
satisfying:
\begin{enumerate}[label=\textup{(\Roman*)}]
\item \textbf{Well-definedness and equivariance.}
 Each $\delta_\Gamma$ is a well-defined bounded linear map between
 the completed endomorphism spaces, equivariant under
 $\operatorname{Aut}(\Gamma)$.
\item \textbf{Coassociativity.}
 For any two composable stable bordered graphs
 $\Gamma_1 \circ_e \Gamma_2$ \textup{(}$\Gamma_2$ inserted at an
 internal edge~$e$ of~$\Gamma_1$\textup{)}, the cocomposition maps
 satisfy
 \begin{equation}\label{eq:heisenberg-coassociativity}
 (\delta_{\Gamma_1} \otimes \id)
 \circ \delta_{\Gamma_2}
 \;=\;
 \delta_{\Gamma_1 \circ_e \Gamma_2}.
 \end{equation}
\item \textbf{Self-sewing = trace.}
 For the non-separating degeneration
 $\Gamma_{\mathrm{ns}} = (\text{vertex of genus } g{-}1
 \text{ with a self-loop})$, the cocomposition is the categorical
 trace:
 \begin{equation}\label{eq:heisenberg-self-sewing}
 \delta_{\Gamma_{\mathrm{ns}}}
 \;=\;
 \operatorname{Tr}_{\Cop}\!\bigl(-\bigr)
 \;\colon\;
 \operatorname{End}(\cF_\lambda)
 \to
 \operatorname{End}(\cF_\lambda).
 \end{equation}
 At genus~$1$, this recovers the annulus identification
 $HH_*(\cH_{\kappa,b}) \simeq \operatorname{Tr}_{\cH_\kappa}$.
\item \textbf{Closed-sector compatibility.}
 The closed projection recovers the genus tower :
 $\pi_{\mathrm{cl}}(\delta_\Gamma(Z_g))
 = \kappa \cdot \lambda_g^{\mathrm{FP}}$ for all~$g$.
\item \textbf{Scalar determinant amplitude.}
 The obstruction-weighted scalar determinant amplitude is
 \begin{equation}\label{eq:heisenberg-cooperad-partition}
 Z_g^{\mathrm{scal}}(\cH_\kappa;\,\Omega)
 \;=\;
 (\det\operatorname{Im}\Omega)^{-\kappa/2}
 \cdot (\zetareg\Delta_{\Sigma_g})^{-\kappa/2}.
 \end{equation}
 For an actual rank-$d$ level-one free-boson theory this specializes
 to $\kappa=d$; for rank one at arbitrary level it is the scalar
 obstruction lane. The unweighted oscillator character remains the
 rank-one Fock character.
\item \textbf{Determination by $\Theta^{\mathrm{oc}}$.}
 The scalar closed projection and the trace-class completed-Fock
 cocompositions are determined by the Heisenberg component of
 $\Theta^{\mathrm{oc}}_{\cH_\kappa}$: the Getzler--Kapranov
 correspondence~\cite{GeK98} identifies these restricted
 cocompositions with evaluations of the MC element on stable bordered
 graphs.
\end{enumerate}
This proves the class-G completed-Fock model of the cooperad axioms.
It does not construct the full modular cooperad on arbitrary boundary
conditions of the global open factorization category.
\end{theorem}

\begin{proof}
The proof exploits the one-particle Bergman reduction
(Theorem~\ref{thm:heisenberg-one-particle-sewing}): every statement
about the Fock-space cooperad reduces to an operator-algebraic
statement on the one-particle Bergman space $A^2(D)$, where the
arguments are elementary.

\medskip\noindent
\textit{Proof of~\textup{(I)}: well-definedness and equivariance.}
Fix a stable bordered graph $\Gamma$ of type $(g,n)$ with internal
edges $E(\Gamma)$. Each internal edge $e$ contributes a sewing
operator $\sewop_{q_e}$ with exponentially decaying matrix
coefficients $|(\sewop_{q_e})_{n,m}| \leq C\,q_e^{(n+m)/2}$
(Lemma~\ref{V2-lem:thqg-X-composition-decay}). In particular,
$\sewop_{q_e}$ is trace class on $A^2(D)$ with
$\|\sewop_{q_e}\|_1 = q_e/(1-q_e)$.

The full cocomposition is the second quantization
$\secquant(\cdot)^\kappa$ applied to a product of trace-class
operators on the one-particle space. By
Proposition~\ref{V2-prop:thqg-X-second-quantization}, $\secquant(T)$ is
trace class on $\operatorname{Sym} A^2(D)$ whenever $T$ is trace
class with $\|T\| < 1$. Since $\|\sewop_{q_e}\| = q_e < 1$, the
composed operator at each vertex is trace class, and $\delta_\Gamma$
is well-defined as a bounded operator.

For equivariance: an automorphism
$\sigma \in \operatorname{Aut}(\Gamma)$ permutes the internal edges
and vertices, and the trace over Fock states is invariant under
permutation of the sewing channels (the trace is cyclically invariant
and the Bergman operators on disjoint edges commute). Hence
$\delta_\Gamma$ descends to the quotient by
$\operatorname{Aut}(\Gamma)$.

\medskip\noindent
\textit{Proof of~\textup{(II)}: coassociativity.}
This is the central structural property. We prove it in three
steps.

\emph{Step~1: reduction to one-particle operators.}
By the one-particle sewing reduction
(Theorem~\ref{thm:heisenberg-one-particle-sewing}(iii)), every
multi-particle sewing amplitude on $\operatorname{Sym} A^2(D)$ is
the second quantization of a one-particle operator:
$A_P = \secquant(R)$ for a pair-of-pants amplitude. The Fredholm
determinant formula
(Theorem~\ref{thm:heisenberg-one-particle-sewing}(iv))
implies that the trace of any composed sewing operator on the Fock
space factors through the one-particle trace:
\begin{equation}\label{eq:heisenberg-one-particle-factorization}
\operatorname{Tr}_{\mathrm{Sym}}\!\bigl(
\secquant(T_1 \circ T_2)\bigr)
\;=\;
\det(1 - T_1 \circ T_2)^{-\kappa}.
\end{equation}
It therefore suffices to establish coassociativity at the one-particle
level.

\emph{Step~2: composition of one-particle operators is associative.}
On $A^2(D)$, the sewing operators $\sewop_{q_e}$ associated to
distinct internal edges $e_1, e_2$ of $\Gamma$ act on disjoint tensor
factors of $A^2(D)^{\otimes |E(\Gamma)|}$ (each edge has its own
pair of half-edges, and the sewing operator connects the Bergman
spaces of the two disks at the node). Operators on disjoint tensor
factors commute:
\begin{equation}\label{eq:heisenberg-operator-commutation}
\sewop_{q_{e_1}} \circ \sewop_{q_{e_2}}
\;=\;
\sewop_{q_{e_2}} \circ \sewop_{q_{e_1}}
\quad\text{on}\;
A^2(D_{e_1}) \otimes A^2(D_{e_2}).
\end{equation}

For composable degenerations $\Gamma_1 \circ_e \Gamma_2$: the
degeneration first along $\Gamma_2$ produces operators $\sewop_{q_f}$
for $f \in E(\Gamma_2)$, and then along $\Gamma_1$ produces operators
$\sewop_{q_e}$ for $e \in E(\Gamma_1)$. The composition acts on
the one-particle space as
\[
\sewker_{\Gamma_1 \circ_e \Gamma_2}
\;=\;
\Bigl(\prod_{e' \in E(\Gamma_1)} \sewop_{q_{e'}}\Bigr)
\circ
\Bigl(\prod_{f \in E(\Gamma_2)} \sewop_{q_f}\Bigr),
\]
and the analogous expression with $\Gamma_1$ inserted into $\Gamma_2$
gives the same operator (since the operators on disjoint edges commute
and the shared edge contributes the same sewing parameter in both
orderings).

\emph{Step~3: Fredholm determinant multiplicativity.}
At the level of partition functions (closed amplitudes), the
coassociativity is witnessed by the multiplicativity of Fredholm
determinants. For disjoint sewing kernels
$\sewker_1, \sewker_2$ on disjoint one-particle Hilbert spaces:
\begin{equation}\label{eq:heisenberg-fredholm-multiplicativity}
\det(1 - \sewker_1 \oplus \sewker_2)
\;=\;
\det(1 - \sewker_1) \cdot \det(1 - \sewker_2).
\end{equation}
For the interaction correction when the sewing kernels share
a Bergman space (composable degenerations), the Fredholm
determinant factorises as
\begin{equation}\label{eq:heisenberg-fredholm-composition}
\det(1 - \sewker_1 \circ \sewker_2)
\;=\;
\det\bigl((1 - \sewker_1)(1 - \sewker_2)
(1 - (1-\sewker_1)^{-1}(1-\sewker_2)^{-1}
\sewker_1\,\sewker_2)\bigr),
\end{equation}
but the key point is that the associativity
$(\sewker_1 \circ \sewker_2) \circ \sewker_3
= \sewker_1 \circ (\sewker_2 \circ \sewker_3)$
holds strictly as an identity of trace-class operators on $A^2(D)$
(operator composition is associative). The Fredholm determinant,
being a spectral invariant of the composite operator, therefore
satisfies coassociativity automatically:
\[
\operatorname{Tr}_{\mathrm{Sym}}\!\bigl(
\delta_{\Gamma_1} \circ \delta_{\Gamma_2}(\cdot)\bigr)
\;=\;
\operatorname{Tr}_{\mathrm{Sym}}\!\bigl(
\delta_{\Gamma_1 \circ_e \Gamma_2}(\cdot)\bigr),
\]
because both sides reduce to the same one-particle operator
$\prod_{e \in E(\Gamma_1 \circ_e \Gamma_2)} \sewop_{q_e}$
on $A^2(D)$, and the second quantization functor $\secquant$ is
a monoidal functor (it preserves composition).

\medskip\noindent
\textit{Proof of~\textup{(III)}: self-sewing = trace.}
The non-separating degeneration $\Gamma_{\mathrm{ns}}$ is the graph
with one vertex of genus $g-1$ and one self-loop. The cocomposition
$\delta_{\Gamma_{\mathrm{ns}}}$ identifies the two half-edges of the
self-loop and traces over the corresponding Fock space:
\[
\delta_{\Gamma_{\mathrm{ns}}}(Z)
\;=\;
\sum_{\mu}
\langle\mu|\, Z\, |\mu\rangle
\;=\;
\operatorname{Tr}_{\cF}(Z),
\]
where $\{|\mu\rangle\}$ is an orthonormal basis of Fock states.

At the one-particle level, the self-loop sewing operator is
$\sewop_q$ on a single copy of $A^2(D)$, and the trace is
$\operatorname{Tr}_{\mathrm{Sym}}(\secquant(\sewop_q)^\kappa)
= \det(1 - \sewop_q)^{-\kappa}$. At genus~$1$
(one pair of pants with a self-loop), this gives
$Z_1 = \det(1 - \sewop_q)^{-\kappa}
= \prod_{n=1}^\infty (1 - q^n)^{-\kappa}
= (q^{-1/24}\,\eta(q))^{-\kappa}$,
reproducing the Heisenberg partition function
(Theorem~\ref{thm:heisenberg-one-particle-sewing}).

This is the categorical trace $\operatorname{Tr}_{\Cop}$ on the
dualizable category $\Cop(\cH_\kappa)$
(Proposition~\ref{prop:heisenberg-open-sector}(iii)): the
non-separating degeneration creates a handle by identifying two
boundary circles, which is precisely the categorical trace
operation $\operatorname{Tr}\colon \operatorname{End}(\cF_\lambda)
\to \operatorname{End}(\cF_\lambda)$. Agreement with the annulus
trace theorem (Theorem~\ref{thm:thqg-annulus-trace}) follows:
\begin{equation}\label{eq:heisenberg-annulus-trace-agreement}
\int_{S^1_p} \Cop(\cH_\kappa)
\;\simeq\;
HH_*(\Cop(\cH_\kappa))
\;\simeq\;
\bigoplus_\lambda \mathbb{C},
\end{equation}
where the categorical Hochschild homology of the semisimple category
$\Cop(\cH_\kappa)$ is concentrated in degree~$0$ (semisimple
$E_1$-algebras have $HH_n = 0$ for $n \geq 1$), with one summand
for each simple object~$\cF_\lambda$.

\medskip\noindent
\textit{Proof of~\textup{(IV)}: closed-sector compatibility.}
The closed scalar projection $\pi_{\mathrm{cl}}$ of the open/closed MC
element (Theorem~\ref{thm:thqg-oc-projection}(i)) sends
$\Theta^{\mathrm{oc}}_{\cH_\kappa} \mapsto
\pi_{\mathrm{scal}}(\Theta_{\cH_\kappa})
= \kappa \cdot \eta \otimes \Lambda$
(Proposition~\ref{prop:heisenberg-gaussian-termination}). The
genus-$g$ component is
$\pi_{\mathrm{scal}}(\Theta_{\cH_\kappa})|_g
= \kappa \cdot \lambda_g^{\mathrm{FP}}
= F_g(\cH_\kappa)$ which is the genus tower of Theorem~D.

The cooperad cocomposition $\delta_\Gamma$ in the closed sector is
the clutching identity
\begin{equation}\label{eq:heisenberg-clutching-closed}
\xi^*_\Gamma(F_g)
\;=\;
\sum_{\substack{g_1 + g_2 = g \\ g_1 + g_2 + |E| = g}}
F_{g_1} \cdot F_{g_2} \cdot \prod_{e \in E(\Gamma)}
(\text{propagator at } e),
\end{equation}
where the propagator at each edge is the sewing operator restricted
to the scalar (degree-zero) channel. For the Heisenberg, the
propagator is $d\log E(z,w)$ (the weight-$1$ Bergman kernel;
Remark~\ref{rem:propagator-weight-universality}), and the scalar
projection extracts $\kappa$ from each propagator factor. The
compatibility amounts to the identity
\[
\pi_{\mathrm{cl}}(\delta_\Gamma(Z_g))
\;=\;
\xi^*_\Gamma(\kappa \cdot \lambda_g^{\mathrm{FP}}),
\]
which holds because the closed projection commutes with the sewing
operations (Theorem~\ref{thm:thqg-mc-forced-consistency}(a,c)).

\medskip\noindent
\textit{Proof of~\textup{(V)}: scalar determinant amplitude.}
The scalar determinant amplitude is obtained by tracing the completed
Fock cooperad over all stable graphs of type $(g,0)$:
\begin{align*}
Z_g^{\mathrm{scal}}(\cH_\kappa;\,\Omega)
&\;=\;
\sum_{\Gamma \in \mathcal{G}_{g,0}}
\frac{1}{|\operatorname{Aut}(\Gamma)|}
\operatorname{Tr}\bigl(\delta_\Gamma(\id)\bigr)
\\[4pt]
&\;=\;
\sum_{\text{pants decompositions}}
\frac{1}{|\operatorname{Aut}|}
\prod_{\alpha=1}^{3g-3}
\det(1 - \sewop_{q_\alpha})^{-\kappa}
\\[4pt]
&\;=\;
(\det\operatorname{Im}\Omega)^{-\kappa/2}
\cdot (\zetareg\Delta_{\Sigma_g})^{-\kappa/2}.
\end{align*}
The first equality is the definition. The second uses the
one-particle reduction: each graph $\Gamma$ specifies a pants
decomposition, and the cocomposition at each vertex is the
second-quantized Bergman operator, whose trace is the Fredholm
determinant (Proposition~\ref{V2-prop:thqg-X-second-quantization}).
The third is the Polyakov--Alvarez formula
(\eqref{eq:thqg-X-polyakov}, Theorem~\ref{V2-thm:thqg-X-heisenberg-sewing-full}(IV)),
which converts the product of one-particle determinants to
the period-matrix determinant times the zeta-regularized
Laplacian determinant. (The summation over pants decompositions
is redundant: all decompositions give the same Fredholm
determinant by
Remark~\ref{rem:thqg-X-pants-independence}.)

\medskip\noindent
\textit{Proof of~\textup{(VI)}: determination by
$\Theta^{\mathrm{oc}}$.}
The Heisenberg is class~G (shadow depth $r_{\max} = 2$;
Proposition~\ref{prop:heisenberg-gaussian-termination}), so the
shadow obstruction tower terminates:
$\pi_{\mathrm{scal}}(\Theta_{\cH_\kappa})
= \kappa \cdot \eta \otimes \Lambda$
with all higher shadow obstructions vanishing ($o_r = 0$ for
$r \geq 3$).
The closed scalar projection is therefore determined by two
operations: separating clutching (binary cocomposition) and
non-separating clutching (genus-raising self-sewing). Both are
one-parameter families controlled by the single scalar~$\kappa$.

The Getzler--Kapranov correspondence~\cite{GeK98} identifies the
cooperad structure maps with evaluations of the MC element
$\Theta^{\mathrm{oc}}_{\cH_\kappa}$ on stable bordered graphs:
the MC equation
\[
D^{\mathrm{oc}}\Theta^{\mathrm{oc}}
+ \tfrac{1}{2}[\Theta^{\mathrm{oc}}, \Theta^{\mathrm{oc}}]
+ \hbar\,\Delta_{\mathrm{clutch}}(\Theta^{\mathrm{oc}})
= 0
\]
(Theorem~\ref{thm:thqg-oc-mc-equation}) encodes all cooperad
relations simultaneously. For the Heisenberg,
$\pi_{\mathrm{scal}}(\Theta^{\mathrm{oc}}_{\cH_\kappa})$ is the
single scalar class $\kappa$, while the open trace-class component is
the second-quantized Bergman sewing operator. Within the completed
Fock packet, the cocompositions $\delta_\Gamma$ of (I)--(V) are
recovered as evaluations
$\delta_\Gamma = \operatorname{ev}_\Gamma(\Theta^{\mathrm{oc}})$.
This is a restricted determination statement, not a construction of
the global modular cooperad on all boundary conditions.
Coassociativity (II) is the chain-level identity encoding the
codimension-$2$ stratum cancellation on
$\overline{\mathcal{M}}_{g,n}$.
\end{proof}

\begin{remark}[Why the Heisenberg is special]
\label{rem:heisenberg-cooperad-special}
\index{Heisenberg!cooperad simplification}
The proof above uses three properties special to the
Heisenberg algebra:
\begin{enumerate}[label=\textup{(\alph*)}]
\item \emph{One-particle reduction.} All multi-particle amplitudes
 factor through one-particle Bergman operators (Wick's theorem).
 This reduces the cooperad from an infinite-dimensional categorical
 structure to finite-dimensional linear algebra on $A^2(D)$.
\item \emph{Semisimplicity.} The open-sector category
 $\Cop(\cH_\kappa)$ is semisimple. There are no higher $A_\infty$
 obstructions, no extension classes, no derived structure beyond
 degree~$0$.
\item \emph{Gaussian termination.} The shadow obstruction tower terminates
 at degree~$2$ (Proposition~\ref{prop:heisenberg-gaussian-termination}),
 so the only MC data entering the cooperad structure is the scalar
 $\kappa$.
\end{enumerate}
For non-Gaussian algebras (classes~L, C, M), the cooperad
construction requires the full open/closed MC element
$\Theta^{\mathrm{oc}}_\cA$ and the higher-degree shadow data
$(C, Q, \ldots)$; the categorical cocomposition maps involve
genuine derived structure (higher $A_\infty$ operations on the
open sector) and the coassociativity proof must use the full MC
equation rather than operator-algebraic arguments.
\end{remark}

\begin{remark}[Trace-class inputs for the Heisenberg cooperad]
\label{rem:ct2-heisenberg-dependencies}
\index{CT-2!Heisenberg!dependencies}
The trace-class cooperad proof has five mathematical inputs.
\begin{enumerate}[label=\textup{(\arabic*)},nosep]
\item the Heisenberg sewing theorem
 (Theorem~\ref{thm:heisenberg-sewing-vol1}): the sewing envelope is
 $\operatorname{Sym} A^2(D)$;
\item the one-particle Bergman reduction
 (Theorem~\ref{thm:heisenberg-one-particle-sewing}): pair-of-pants
 sewing is $\secquant(R)$, and every closed amplitude is
 $\det(1 - \sewker)^{-\kappa}$;
\item the bordered FM compactification
 (Construction~\ref{constr:bordered-fm}): codimension-$2$
 corner cancellation gives $\partial^2 = 0$;
\item the annulus trace theorem
 (Theorem~\ref{thm:thqg-annulus-trace}):
 $\int_{S^1_p} \Cop \simeq HH_*(\Cop)$;
\item the open/closed MC equation
 (Theorem~\ref{thm:thqg-oc-mc-equation}):
 $(D^{\mathrm{oc}})^2 = 0$ and
 $\Theta^{\mathrm{oc}}_\cA \in \operatorname{MC}(\mathfrak{g}^{\mathrm{oc}})$
 at all genera.
\end{enumerate}
Together they identify the Fredholm-determinant amplitudes on the
completed Fock open-sector packet with the cocompositions of a
trace-class cooperad. The verification is operator-algebraic:
associativity of composition on $A^2(D)$ and multiplicativity of
Fredholm determinants give the coassociativity identities on the
completed-Fock class-G packet.
\end{remark}

\begin{remark}[Cooperad vs Feynman transform]
\label{rem:heisenberg-cooperad-feynman}
\index{Feynman transform!cooperad relation}
The modular cooperad on $\Cop(\cH_\kappa)$ is dual to the
modular operad structure on the genus-graded bar complex
$\barB^{\mathrm{full}}(\cH_\kappa)$
(Theorem~\ref{thm:chain-modular-functor}). The bar complex
carries the structure of an algebra over the Feynman transform
$\operatorname{FT}_{\mathcal{M}\mathrm{od}}$
of the Getzler--Kapranov modular operad
\cite{GeK98}; the cooperad on the open sector is the dual
structure obtained by taking categorical traces. The
coassociativity of the cooperad
(Theorem~\ref{thm:heisenberg-modular-cooperad}(II))
is equivalent to the associativity of the modular operad
composition on the bar side (the codimension-$2$ stratum identity
on~$\overline{\mathcal{M}}_{g,n}$).
\end{remark}

\section{\texorpdfstring{$\Delta_5$}{Delta5} normalization firewall}
\label{sec:heisens-supplement}

\begin{remark}[Heisenberg boundary and the $\Delta_5$ normalization firewall]
\label{rem:heisens-1}
The class-G Heisenberg boundary at the Borel cusp of
$\bar{\mathcal A}_2^{\mathrm{BB}}$ records the Gaussian,
semi-abelian degeneration of the K3 Mukai-lattice sector. It does
not produce $\Delta_5$, and it does not reconstruct the chiral
bialgebra $\mathbf H_{\Delta_5}$.

The automorphic input is separate. In Gritsenko's normalization, the
Igusa cusp form of weight $5$ is obtained from the Jacobi cusp form
$\eta^9\vartheta_1$ of weight $5$ and index $1/2$:
\[
\Delta_5=\operatorname{Grit}(\eta^9\vartheta_1)
\qquad
\text{\cite{Gritsenko1999,EichlerZagier1985}.}
\]
The Borcherds-product normalization realizes the same automorphic
divisor from the weak Jacobi form $\phi_{0,1}^{K3}$, with the usual
normalization constant separated
\cite{Borcherds1998,GritsenkoNikulin1998}. The
Oberdieck--Pandharipande reduced-DT convention uses
$D_5=64^{-1}\Delta_5$; this is a normalization bridge, not a
Heisenberg $\kappa$-value. At the Mukai Heisenberg boundary the
chiral scalar is $\kappa_{\mathrm{ch}}=4$ in the Route-B convention,
whereas the Borcherds denominator has weight $5$. These are distinct
invariants.

Any Andrianov-style spinor $L$-factorization for a normalized Siegel
eigenline belongs to the automorphic
$\Delta_5$ problem \cite{Andrianov1974,Arthur2013Endoscopic}. It is
not a consequence of the Gaussian Heisenberg boundary data proved in
this section.
\end{remark}
